%===============================================================================
%  DISCARD ON MERGE --- mirrors book/main.tex
%
%  Nothing between this line and the marker "END OF DISCARD BLOCK" belongs in
%  [B].  The macros \projtens, \convol, \projmodtens, \isom are copied verbatim
%  from book/main.tex; the operators sep, sepd, asep, asepd, sng, dist and the
%  four constants of the algebra are the preamble additions that Part IV asks
%  to be vetted.  The body below uses only [B]'s macros.
%
%  Section and theorem numbers in this file are an artefact of the standalone
%  document: everything is cited by label, never by number.
%===============================================================================
\documentclass[11pt]{article}
\usepackage[margin=1.05in]{geometry}
\usepackage[T1]{fontenc}
\usepackage{amsmath,amssymb,amsthm}
\usepackage{mathrsfs}
\usepackage{enumitem}

\theoremstyle{plain}
\newtheorem{theorem}{Theorem}[section]
\newtheorem{lemma}[theorem]{Lemma}
\newtheorem{corollary}[theorem]{Corollary}
\newtheorem{proposition}[theorem]{Proposition}
\theoremstyle{definition}
\newtheorem{definition}[theorem]{Definition}

% --- copied verbatim from book/main.tex -------------------------------------
\newcommand{\projtens}{\mathbin{\widehat{\otimes}}}
\newcommand{\convol}{\ast}
\newcommand{\projmodtens}[1]{\mathbin{\widehat{\otimes}}_{#1}}
\newcommand{\isom}[1]{\mathrel{\mathop{\cong}\limits_{#1}}}

% --- preamble additions required in book/main.tex (see Part IV) --------------
\DeclareMathOperator{\dist}{dist}
\DeclareMathOperator{\sep}{sep}
\DeclareMathOperator{\sepd}{sepd}
\DeclareMathOperator{\asep}{asep}
\DeclareMathOperator{\asepd}{asepd}
\DeclareMathOperator{\sng}{sng}
\newcommand{\cri}{c_{\mathrm{ri}}}
\newcommand{\cli}{c_{\mathrm{li}}}
\newcommand{\crai}{c_{\mathrm{rai}}}
\newcommand{\clai}{c_{\mathrm{lai}}}

% --- proofreading convenience -----------------------------------------------
% The labels below live in [B], not here; the stubs make every reference to
% them print as the label itself, so that the file compiles clean on its own.
% Delete the whole block on merge, where the labels resolve for real.
\makeatletter
\DeclareRobustCommand{\cite}[2][]{[\texttt{#2}\ifx\\#1\\\else,~#1\fi]}
\makeatother
\newlabel{AnnihCModIsRetAnnihMod}{{\texttt{AnnihCModIsRetAnnihMod}}{--}}
\newlabel{CStarAlgIsTopFlatOverItsIdeal}{{\texttt{CStarAlgIsTopFlatOverItsIdeal}}{--}}
\newlabel{DualMetTopProjIsMetrInj}{{\texttt{DualMetTopProjIsMetrInj}}{--}}
\newlabel{GoodIdealMetTopProjIsUnital}{{\texttt{GoodIdealMetTopProjIsUnital}}{--}}
\newlabel{L1MetTopProjAndMetrFlatOfMeasAlg}{{\texttt{L1MetTopProjAndMetrFlatOfMeasAlg}}{--}}
\newlabel{MeasAlgIsL1TopFlat}{{\texttt{MeasAlgIsL1TopFlat}}{--}}
\newlabel{MetCTopFlatMod}{{\texttt{MetCTopFlatMod}}{--}}
\newlabel{MetCTopFreeMod}{{\texttt{MetCTopFreeMod}}{--}}
\newlabel{MetCTopInjMod}{{\texttt{MetCTopInjMod}}{--}}
\newlabel{MetCTopProjMod}{{\texttt{MetCTopProjMod}}{--}}
\newlabel{MetCTopProjModViaCanonicMorph}{{\texttt{MetCTopProjModViaCanonicMorph}}{--}}
\newlabel{MetInjIsOneTopInj}{{\texttt{MetInjIsOneTopInj}}{--}}
\newlabel{MetTopFlatAnnihModCharac}{{\texttt{MetTopFlatAnnihModCharac}}{--}}
\newlabel{MetTopFlatCycModCharac}{{\texttt{MetTopFlatCycModCharac}}{--}}
\newlabel{MetTopFlatModCoProd}{{\texttt{MetTopFlatModCoProd}}{--}}
\newlabel{MetTopFlatModEssQuotIsFlatBanSp}{{\texttt{MetTopFlatModEssQuotIsFlatBanSp}}{--}}
\newlabel{MetTopInjModAnnihPartIsInjBanSp}{{\texttt{MetTopInjModAnnihPartIsInjBanSp}}{--}}
\newlabel{MetTopInjModCCharac}{{\texttt{MetTopInjModCCharac}}{--}}
\newlabel{MetTopInjModProd}{{\texttt{MetTopInjModProd}}{--}}
\newlabel{MetTopInjOfAnnihModCharac}{{\texttt{MetTopInjOfAnnihModCharac}}{--}}
\newlabel{MetTopProjCycModCharac}{{\texttt{MetTopProjCycModCharac}}{--}}
\newlabel{MetTopProjIsMetTopFlat}{{\texttt{MetTopProjIsMetTopFlat}}{--}}
\newlabel{MetTopProjModCCharac}{{\texttt{MetTopProjModCCharac}}{--}}
\newlabel{MetTopProjModEssQuotIsProjBanSp}{{\texttt{MetTopProjModEssQuotIsProjBanSp}}{--}}
\newlabel{MetTopProjOfAnnihModCharac}{{\texttt{MetTopProjOfAnnihModCharac}}{--}}
\newlabel{MetTopRelFlatModCompIdealPartCompl}{{\texttt{MetTopRelFlatModCompIdealPartCompl}}{--}}
\newlabel{MetTopRelFlatModPartCompl}{{\texttt{MetTopRelFlatModPartCompl}}{--}}
\newlabel{MetTopRelInjModAnnihCompl}{{\texttt{MetTopRelInjModAnnihCompl}}{--}}
\newlabel{MetTopRelInjModCompIdealAnnihCompl}{{\texttt{MetTopRelInjModCompIdealAnnihCompl}}{--}}
\newlabel{MetTopRelProjModCompIdealPartCompl}{{\texttt{MetTopRelProjModCompIdealPartCompl}}{--}}
\newlabel{MetTopRelProjModPartCompl}{{\texttt{MetTopRelProjModPartCompl}}{--}}
\newlabel{SectionPreliminariesOnHarmonicAnalysis}{{\texttt{SectionPreliminariesOnHarmonicAnalysis}}{--}}
\newlabel{SubSectionBanachAlgebrasAndTheirModules}{{\texttt{SubSectionBanachAlgebrasAndTheirModules}}{--}}
\newlabel{SubSectionBanachSpaces}{{\texttt{SubSectionBanachSpaces}}{--}}
\newlabel{SubSectionC0SModules}{{\texttt{SubSectionC0SModules}}{--}}
\newlabel{SubSectionL1GModules}{{\texttt{SubSectionL1GModules}}{--}}
\newlabel{SubSectionMetricAndTopologicalFlatness}{{\texttt{SubSectionMetricAndTopologicalFlatness}}{--}}
\newlabel{SubSectionc0AndlInftyModules}{{\texttt{SubSectionc0AndlInftyModules}}{--}}
\newlabel{c0AndlInftyModlIfty}{{\texttt{c0AndlInftyModlIfty}}{--}}

\title{The essential and the annihilator part:\\
material for section 2.2 of \emph{Metric and topological homology},\\
with the accompanying insertions in chapters 1 and 3}
\author{Book-ready form of the companion note}
\date{25 August 2026 \\[2pt] {\small version \texttt{20260825\_0829};
supersedes \texttt{20260825\_0530}}}
%===============================================================================
%  END OF DISCARD BLOCK
%===============================================================================

\begin{document}
\maketitle

\noindent
This file is the companion note in the form in which it drops into [B]. It has
four parts. Part I holds the prose paragraphs that belong in chapter 1, each
headed by where it goes. Part II is the new section 2.2, ``The essential and the
annihilator part'', in six subsections; at the head of each subsection the
statements of [B] that are to be interleaved are listed by label, with the point
of insertion. Part III holds the material for chapter 3. Part IV is the merge
manifest and the list of preamble additions.

Throughout, $\sep_X(x)=\dist(x,X_{ess})$ is the \emph{separation} of a vector
from the essential part and
$\sepd(X)=\sup\{\sep_X(x):x\in B_{X_{ann}}\}$ its \emph{separation defect};
replacing annihilator vectors by nets annihilated only asymptotically gives a
second defect $\asepd(X)$ and the chain $\sepd(X)\leq\sepd(X^*)\leq\asepd(X)$.
Two constants of the algebra arise by evaluating the defect on the free module
of rank one and on its dual,
\[
\sepd(A_+)=\frac{1}{1+\cri(A)},
\qquad
\sepd(A_+^*)=\frac{1}{1+\crai(A)},
\]
and the three homological properties then obey one estimate with two constants:
a $\langle$~metrically / $C$-topologically~$\rangle$ projective $P$, injective
$J$ or flat $F$ satisfies
$\sepd(P)\leq\langle\,\sepd(A_+)\ /\ C\sepd(A_+)\,\rangle$ and
$\sepd(J),\sepd(F)\leq\langle\,\sepd(A_+^*)\ /\ C\sepd(A_+^*)\,\rangle$.
Since $\sepd=1$ on a non-zero annihilator module, these lines return the algebra
clauses of the annihilator characterisations collected in subsection 2.2.6 --- a
right identity in the projective case, a bounded right approximate identity in
the other two.

\section{For chapter 1}\label{PartI}

The paragraphs below are prose: chapter 1 of [B] carries no numbered statement
anywhere, and these do not change that. Each is headed by its destination.

\subsection{For subsection~\ref{SubSectionBanachSpaces}, in the paragraph that
already carries $E^*/F^\perp\isom{\mathbf{Ban}_1}F^*$ and
${(E/F)}^*\isom{\mathbf{Ban}_1}F^\perp$}

For a non-empty subset $M$ of a Banach space $E$ and a vector $x\in E$ we write
$\dist(x,M)=\inf \{\lVert x-y\rVert:y\in M \}$. The function
$\dist(\cdot,M):E\to\mathbb{R}_+$ is $1$-Lipschitz, and it is positively
homogeneous whenever $M$ is a linear subspace. If $M$ is a closed subspace, the
Hahn-Banach theorem gives the dual description
\[
\dist(x,M)=\max \{|f(x)|:f\in B_{M^\perp} \},
\]
the maximum being attained; combined with the isometric isomorphism
$i:(E/M)^* \to M^\perp$ above, this says that
\[
\dist(x,M)=\lVert x+M\rVert_{E/M}=\max \{ |f(x)| : f\in B_{{(E/M)}^*} \}
\]
We shall use three further classical facts about weak topologies. By the 
Banach-Alaoglu theorem the closed unit ball of a dual space is weak${}^*$ 
compact. By Goldstine's theorem
$\iota_E(B_E)$ is weak${}^*$ dense in $B_{E^{**}}$. By Mazur's theorem the weak
closure of a convex subset of $E$ coincides with its norm closure.

\subsection{For subsection~\ref{SubSectionBanachAlgebrasAndTheirModules}, after
the paragraph introducing approximate identities}

It is convenient to name the optimal bounds. By
$\langle$~$\cri(A)$ / $\cli(A)$~$\rangle$ we denote the infimum of the norms of
the $\langle$~right / left~$\rangle$ identities of $A$, and by
$\langle$~$\crai(A)$ / $\clai(A)$~$\rangle$ the infimum of the bounds of the
$\langle$~right / left~$\rangle$ approximate identities of $A$; each of the four
is $+\infty$ if there is nothing to take the infimum over, and each of the four
is $0$ for the zero algebra, whose only element is a two-sided identity and
whose only net is a two-sided approximate identity. For $A\neq \{0 \}$ all four
constants lie in $[1,+\infty]$. If $A$ is commutative, a left
approximate identity is a right one and conversely, so $\clai(A)=\crai(A)$; it is
not known whether $\clai(A)=\crai(A)$ whenever both are finite. If $A$ is unital
then
\[
\cri(A)=\cli(A)=\crai(A)=\clai(A)=\lVert e_A\rVert ,
\]
the zero algebra included, with $e_A=0$ and all four constants $0$. Finally, 
when $\langle$~$\clai(A)$ / $\crai(A)$~$\rangle$ is finite it is attained: some 
 $\langle$~left / right~$\rangle$approximate identity is 
bounded by $\langle$~$\clai(A)$ / $\crai(A)$~$\rangle$ itself. This is 
theorem~\ref{SepOfDualUnitization} below.

\subsection{After introduction of essential and faithful modules add the 
following facts}.

Suppose $A$ has a $\langle$~left / right~$\rangle$ approximate identity 
${(e_\nu)}_{\nu\in N}$. For any $\langle$~left / right~$\rangle$ Banach $A$-module 
$X$ and any $x\in X_{ess}$ the net $\langle$~${(e_\nu \cdot x)}_{\nu\in N}$ / 
${(x\cdot e_\nu)}_{\nu\in N}$~$\rangle$ converges to $x$. If $Y$ is another 
$\langle$~left / right~$\rangle$ Banach $A$-module and $\phi:X\to Y$ is an 
$A$-mprphism, then $\phi(X_{ess})\subset Y_{ess}$ 
and $\phi(X_{ann})\subset Y_{ann}$. 

\subsection{Move this to the new subsection 2.2 before proposition 2.4 of this 
note as a numbered definition}

An \emph{annihilator net} in a $\langle$~left / right~$\rangle$ Banach
$A$-module $X$ is a net ${(x_\nu)}_{\nu\in N}$ in $X$ such that
$\langle$~${(a\cdot x_\nu)}_{\nu\in N}$ / ${(x_\nu\cdot a)}_{\nu\in N}$~$\rangle$ 
is norm null for every $a\in A$.

\subsection{For subsection~\ref{SubSectionBanachAlgebrasAndTheirModules}, in the
paragraph that already records
${(X^*)}^{\perp S}={\operatorname{cl}_X(S X)}^\perp$}

Setting $S=A$ in the identity above and reading it on the two sides at once: for
a left Banach $A$-module $X$ the dual right module $X^*$ has
${(X^*)}_{ann}={(X_{ess})}^\perp$. The two subspaces ${(X^*)}_{ess}$ 
and ${(X_{ann})}^\perp$ are in general different, though the first is always 
contained in the second.

\section{The essential and the annihilator part (the new section 2.2)}\label{PartII}

The subsection headings below are those of the target structure. A new section
2.2 is inserted between the present sections 2.1 and 2.2 of [B]; the present
section 2.2 becomes 2.3 and the present 2.3 becomes 2.4, save for the present
subsection 2.2.1, which becomes subsection 2.2.6 below.

\subsection{Separation}\label{SubSectionSeparation}

\emph{Nothing of [B] is interleaved here.}

\begin{definition}\label{SepOfMod}
Let $X$ be a $\langle$~left / right~$\rangle$ Banach $A$-module. 
The \emph{separation} of a vector
$x\in X$ is
\[
\sep_X(x):=\dist(x,X_{ess}),
\]
and the \emph{separation defect} of $X$ is
\[
\sepd(X):=\sup \{\sep_X(x):x\in B_{X_{ann}} \}.
\]
For $x\in X$ we write
$\sng(x):=\inf \{\lVert a\cdot x\rVert/\lVert a\rVert:a\in A\setminus \{0 \} \}$
for the \emph{singularity} of $x$, so for $x\in X_{ann}$ we have
$\sng(x)=0$.
\end{definition}

The essential part is a closed linear subspace, so $\sep_X$ is $1$-Lipschitz and
positively homogeneous; we use both properties below without further comment. In
particular
$\sepd(X)=\sup \{\sep_X(x)/\lVert x\rVert:x\in X_{ann}\setminus \{0 \} \}$ as
soon as $X_{ann}\neq \{0 \}$, and $\sepd(X)=\sep_X(0)=0$ otherwise. The defect
measures the directed gap from the annihilator part to the essential part. The 
degenerate case is not idle: modules with $X_{ann}= \{0 \}$ but $X_{ess}\neq X$ 
occur constantly and for them $\sepd(X)=0$ while the module is far from 
essential.

\begin{lemma}\label{SepLowerBound}
Suppose $A$ has a $\langle$~left / right~$\rangle$ approximate identity bounded
by $C$. Then in every $\langle$~left / right~$\rangle$ Banach $A$-module $X$, 
for any $x\in X_{ann}$ holds
\[
\sep_X(x)\geq\frac{\lVert x\rVert}{1+C}\qquad(x\in X_{ann}).
\]
Note: one can put $\langle$~$C=\clai(A)$ / $C=\crai(A)$~$\rangle$.
\end{lemma}

\begin{proof} 
We write the left-hand case, the right-hand one is a trivial adaptation. 
Let ${(e_\nu)}_{\nu\in N}$ be a left approximate identity of $A$ 
with $\sup_\nu\lVert e_\nu\rVert\leq C$. Dispose 
first of $C=0$: then $A=\{0\}$, $X_{ess}=0$. So the assertion reads holds with 
equality. Assume from now on $C>0$.

Now let $x\in X_{ann}$ and $y\in X_{ess}$. Since $e_\nu\cdot x=0$, then
\[
\lVert e_\nu\cdot y\rVert=\lVert e_\nu\cdot(y-x)\rVert
\leq\lVert e_\nu\rVert\lVert x-y\rVert\leq K\lVert x-y\rVert ,
\]
and letting $\nu$ run through $N$ turns the left-hand side into $\lVert g\rVert$. 
So $\lVert x-y\rVert\geq\lVert y\rVert/C$, while the triangle inequality gives
$\lVert x-y\rVert\geq\lVert x\rVert-\lVert y\rVert$. If
$\lVert y\rVert\geq C\lVert x\rVert/(1+C)$ the first estimate gives
$\lVert x-y\rVert\geq\lVert x\rVert/(1+C)$; if
$\lVert y\rVert<K\lVert x\rVert/(1+C)$ the second gives the same. Taking the
infimum over $y\in X_{ess}$ yields $\sep_X(x)\geq\lVert x\rVert/(1+C)$. By taking 
supremum over all upper norm bounds of all left approximate identities 
one can set $\langle$~$C=\clai(A)$ / $C=\crai(A)$~$\rangle$.
\end{proof}

\begin{proposition}\label{SepVsSepDual}
For every $\langle$~left / right~$\rangle$ Banach $A$-module $X$ one has 
$\sepd(X)\leq\sepd(X^*)$.
\end{proposition}

\begin{proof}
If $X_{ann}= \{0 \}$ then $\sepd(X)=0$ and there is nothing to prove. Assume
$X_{ann}\neq \{0 \}$. By the Hahn-Banach description of the distance to a closed
subspace and the identity ${(X^*)}_{ann}={(X_{ess})}^\perp$, we get
\[
\sep_X(x)
=\dist(x, X_{ess})
=\max \{|f(x)|:f\in B_{(X_{ess})^{\perp}} \}
=\max \{|f(x)|:f\in B_{{(X^*)}_{ann}} \}
\]
for all $x\in X$. Therefore,
\[
\sepd(X)
=\sup \{ \sep_X(x) : x\in X_{ann} \}
=\sup \{ |f(x)| : x\in X_{ann}, f\in B_{{(X^*)}_{ann}} \}
\]
On the other hand $(X_{ess})^*\subset(X_{ann})^\perp$. So 
for $f\in{(X^*)}_{ann}\setminus \{0 \}$ holds
\[
\sep_{X^*}(f)=\dist \bigl(f,{(X^*)}_{ess} \bigr)
\geq\dist \bigl(f,{(X_{ann})}^\perp \bigr)
=\Vert f+{(X_{ann})}^{\perp}\Vert_{X/(X_{ann})^\perp}
=\bigl\lVert f|_{X_{ann}}\bigr\rVert ,
\]
the last equality being the isometric identification of $X^*/Y^\perp$ with $Y^*$
by restriction, applied to the closed subspace $Y=X_{ann}$. Taking the supremum 
over $f\in B_{{(X^*)}_{ann}}$ and comparing with the previous
calculation gives 
$
\sepd(X^*)
\geq\sup \bigl \{ \lVert f|_{X_{ann}}\rVert : f\in B_{{(X^*)}_{ann}} \bigr \}
=\sepd(X).
$
\end{proof}

The gap between ${(X^*)}_{ess}$ and ${(X_{ann})}^\perp$ is exactly the slack in
proposition~\ref{SepVsSepDual}. Since ${(X^*)}_{ann}={(X_{ess})}^\perp$, the
defect of the dual is
$\sepd(X^*)=\sup \{\sep_{X^*}(f):f\in{(X_{ess})}^\perp,\,\lVert f\rVert\leq 1 \}$,
and the next proposition computes each of these numbers out of nets in $X$ that
are annihilated only asymptotically. That is what makes it computable.

\begin{proposition}\label{SepOfDualViaAnnihNets}
Let $X$ be a $\langle$~left / right~$\rangle$ $A$-module 
and $f\in{(X_{ess})}^\perp$. Then
\begin{enumerate}[label = (\roman*)]
\item $\sep_{X^*}(f)=\max \{|q(f)| : q\in B_{{({(X^*)}_{ess})}^{\perp}} \}$;

\item $\sep_{X^*}(f)=\sup \bigl \{
  \textstyle\limsup_\nu|f(x_\nu)|:
  {(x_\nu)}_{\nu\in N}\ \text{an annihilator net in }B_X
 \bigr \} $;

\end{enumerate}
\end{proposition}

\begin{proof}
$(i)$ This is a direct conssequence of the formula for distance to ${(X^*)}_{ess}$ 
in terms of functionals vanishing on ${(X^*)}_{ess}$. 

===================================================

$(ii)$ Writing  $L_a:X\to X,\,x\mapsto a\cdot x$, a bounded operator 
with $L_a^*g=g\cdot a$, we have $E=\operatorname{cl}_{X^*}(\operatorname{span} \{L_a^*g \})$, and since
$\langle L_a^{**}q,g\rangle=\langle q,L_a^*g\rangle$ the condition $q|_E=0$
holds if and only if $L_a^{**}q=0$ for every $a\in A$.

Nets give at most the distance. Let ${(x_\nu)}_{\nu\in N}$ be an annihilator net
in $B_X$ and put $L=\limsup_\nu|f(x_\nu)|$. Refining twice --- once so that
$|f(x_\nu)|\to L$, once by Banach-Alaoglu so that $\iota_X(x_\nu)\to q$
weak${}^*$ for some $q\in B_{X^{**}}$ --- and using that a subnet of an
annihilator net is one, we get
$\langle L_a^{**}q,g\rangle=\lim_\nu g(a\cdot x_\nu)=0$ for all $a\in A$ and
$g\in X^*$, because $\lVert a\cdot x_\nu\rVert\to 0$. So $q$ is admissible and
$|q(f)|=L$.

The distance is attained along a net. Let $q$ attain the maximum; after
multiplying by a unimodular scalar we may assume $q(f)=\sep_{X^*}(f)$. By
Goldstine's theorem there is a net ${(y_\mu)}$ in $B_X$ with
$\iota_X(y_\mu)\to q$ weak${}^*$; then $f(y_\mu)\to q(f)$ and, for all $a\in A$
and $g\in X^*$, $g(a\cdot y_\mu)=(g\cdot a)(y_\mu)\to q(g\cdot a)=0$ because
$g\cdot a\in E$. That is, $a\cdot y_\mu\to 0$ weakly for every $a\in A$. Fix a
finite $F\subset A$ and $\epsilon>0$; the set
$ \{ \bigl(f(x),{(a\cdot x)}_{a\in F} \bigr):x\in B_X \}
\subset\mathbb{C}\times X^F$ is convex, being the image of $B_X$ under a linear
map, and $\bigl(q(f),0\bigr)$ lies in its weak closure, hence by Mazur's theorem
in its norm closure; so there is $x\in B_X$ with $|f(x)-q(f)|<\epsilon$ and
$\lVert a\cdot x\rVert<\epsilon$ for every $a\in F$. Directing the pairs
$(F,\epsilon)$ by $(F,\epsilon)\preceq(F',\epsilon')$ when $F\subset F'$ and
$\epsilon'\leq\epsilon$, and choosing such a vector at each index, produces an
annihilator net in $B_X$ along which $|f(\cdot)|$ tends to
$|q(f)|=\sep_{X^*}(f)$.
\end{proof}

\begin{definition}\label{AsympSepOfMod}
The \emph{asymptotic separation} of an annihilator net ${(x_\nu)}_{\nu\in N}$ in
a left $A$-module $X$ is
$\asep_X \bigl({(x_\nu)}_{\nu\in N} \bigr):=\limsup_\nu\sep_X(x_\nu)$, and the
\emph{asymptotic separation defect} of $X$ is
\[
\asepd(X):=\sup \bigl\{
  \asep_X \bigl({(x_\nu)}_{\nu\in N} \bigr):
  {(x_\nu)}_{\nu\in N}\ \text{an annihilator net in }B_X
 \bigr\}\in[0,1].
\]
\end{definition}

\begin{proposition}\label{TheChain}
For every left Banach $A$-module $X$ one has
$\sepd(X)\leq\sepd(X^*)\leq\asepd(X)$.
\end{proposition}

\begin{proof}
The first inequality is proposition~\ref{SepVsSepDual}. For the second, let
$f\in{(X_{ess})}^\perp$ with $\lVert f\rVert\leq 1$; if $0$ is the only such $f$
there is nothing to prove. For $x\in X$ and $g\in X_{ess}$ one has $f(x)=f(x-g)$,
so $|f(x)|\leq\lVert x-g\rVert$, and taking the infimum over $g$ gives
$|f(x)|\leq\sep_X(x)$. Hence
$\limsup_\nu|f(x_\nu)|\leq\limsup_\nu\sep_X(x_\nu)\leq\asepd(X)$ for every
annihilator net in $B_X$, and proposition~\ref{SepOfDualViaAnnihNets} computes
$\sep_{X^*}(f)$ as the supremum of the left-hand side over such nets. Taking the
supremum over $f\in B_{{(X^*)}_{ann}}$ gives $\sepd(X^*)\leq\asepd(X)$.
\end{proof}

The three terms ask the same question with increasing generosity: exact
annihilator vectors, annihilator directions in the dual, asymptotically
annihilated nets. Both inequalities can be strict, and the next theorem says when
they all collapse.

\begin{theorem}\label{AsympSepDirectionalUnital}
Let $X$ be a left Banach $A$-module such that every annihilator net ${(x_\nu)}$
in $B_X$ satisfies $\dist(x_\nu,X_{ann})\to 0$; this holds for every module when
$A$ is unital. Then $\sepd(X)=\sepd(X^*)=\asepd(X)$.
\end{theorem}

\begin{proof}
By proposition~\ref{TheChain} it suffices to prove $\asepd(X)\leq\sepd(X)$. Let
${(x_\nu)}_{\nu\in N}$ be an annihilator net in $B_X$ and $\epsilon>0$. By
hypothesis there is $\nu_0$ such that for $\nu\geq\nu_0$ one may choose
$y_\nu\in X_{ann}$ with $\lVert x_\nu-y_\nu\rVert<\epsilon$, and then
$\lVert y_\nu\rVert\leq 1+\epsilon$; homogeneity of $\sep_X$ and the definition
of $\sepd$ give $\sep_X(y_\nu)\leq(1+\epsilon)\sepd(X)$, trivially so when
$y_\nu=0$. The Lipschitz property then gives
$\sep_X(x_\nu)\leq(1+\epsilon)\sepd(X)+\epsilon$ for $\nu\geq\nu_0$. Take the
limit superior in $\nu$, the supremum over annihilator nets in $B_X$, and let
$\epsilon\to 0$.

Finally, let $A$ be unital and ${(x_\nu)}$ any annihilator net in $B_X$. Put
$y_\nu:=x_\nu-e_A\cdot x_\nu$; then
$a\cdot y_\nu=a\cdot x_\nu-(ae_A)\cdot x_\nu=0$ for every $a\in A$, so
$y_\nu\in X_{ann}$, while
$\lVert x_\nu-y_\nu\rVert=\lVert e_A\cdot x_\nu\rVert\to 0$.
\end{proof}

Both defects are computed on $\bigoplus_1$- and $\bigoplus_\infty$-sums by two
facts: the annihilator part of a sum is the sum of the annihilator parts, and the
distance to a coordinatewise closed subspace is the sum, respectively the
supremum, of the coordinate distances. The essential part is the sum of the
essential parts only in the first case. The lower bounds all come from one
observation, which we isolate because it is used again for free, cofree and
injective modules.

\begin{lemma}\label{SepUnderRetracts}
Let $\pi:X\to Y$ and $\sigma:Y\to X$ be contractive $A$-morphisms with
$\pi\sigma=1_Y$. Then $\sigma(Y_{ann})\subset X_{ann}$,
$\pi(X_{ess})\subset Y_{ess}$ and $\sep_X(\sigma(y))\geq\sep_Y(y)$ for every
$y\in Y$; consequently $\sepd(X)\geq\sepd(Y)$ and $\asepd(X)\geq\asepd(Y)$.
\end{lemma}

\begin{proof}
For $y\in Y_{ann}$ one has $a\cdot\sigma(y)=\sigma(a\cdot y)=0$, and
$\pi(a\cdot x)=a\cdot\pi(x)\in\operatorname{span}(AY)$ while $Y_{ess}$ is closed.
Hence for $y\in Y$ and $z\in X_{ess}$,
\[
\lVert\sigma(y)-z\rVert\geq\lVert\pi\sigma(y)-\pi(z)\rVert
=\lVert y-\pi(z)\rVert\geq\sep_Y(y) ,
\]
and taking the infimum over $z$ gives $\sep_X(\sigma(y))\geq\sep_Y(y)$. Since
also $\lVert\sigma(y)\rVert\leq\lVert y\rVert$, restricting to $y\in Y_{ann}$ and
forming the quotient of the two gives $\sepd(X)\geq\sepd(Y)$. Finally, if
${(y_\nu)}$ is an annihilator net in $B_Y$ then so is ${(\sigma(y_\nu))}$ in
$B_X$, and $\limsup_\nu\sep_X(\sigma(y_\nu))\geq\limsup_\nu\sep_Y(y_\nu)$; hence
$\asepd(X)\geq\asepd(Y)$.
\end{proof}

\begin{proposition}\label{SepOfSums}
Let ${(X_\lambda)}_{\lambda\in\Lambda}$ be left Banach $A$-modules with
$\Lambda\neq\varnothing$, and let $X=\bigoplus_p \{X_\lambda:\lambda\in\Lambda \}$
with the coordinatewise action, for $p=1$ or $p=\infty$. Then
$X_{ann}=\bigoplus_p \{{(X_\lambda)}_{ann}:\lambda\in\Lambda \}$ for both $p$,
and
\begin{enumerate}[label = (\roman*)]
\item $\sepd \bigl(\bigoplus_1 \{X_\lambda:\lambda\in\Lambda \} \bigr)
=\sup_\lambda\sepd(X_\lambda)$;

\item $\sepd \bigl(\bigoplus_\infty \{X_\lambda:\lambda\in\Lambda \} \bigr)
\geq\sup_\lambda\sepd(X_\lambda)$, with equality when $\Lambda$ is finite;

\item $\asepd \bigl(\bigoplus_p \{X_\lambda:\lambda\in\Lambda \} \bigr)
\geq\sup_\lambda\asepd(X_\lambda)$ for $p=1$ and for $p=\infty$, with equality
when $\Lambda$ is finite;

\item for $p=1$,
$X_{ess}=\bigoplus_1 \{{(X_\lambda)}_{ess}:\lambda\in\Lambda \}$ and
$\sep_X(x)=\sum_\lambda\sep_{X_\lambda}(x_\lambda)$ for every $x\in X$.
\end{enumerate}
\end{proposition}

\begin{proof}
The action is contractive for both $p$, and $a\cdot x=0$ for every $a$ if and
only if $a\cdot x_\lambda=0$ for every $a$ and every $\lambda$, which is the
description of $X_{ann}$. Each $X_\mu$ is a $1$-retract of $X$ via the coordinate
injection and projection, so lemma~\ref{SepUnderRetracts} gives all three lower
bounds.

$(iv)$ The inclusion $X_{ess}\subset\bigoplus_1 \{{(X_\lambda)}_{ess} \}$ holds
because $a\cdot x$ lies in the right-hand side, a closed subspace. Conversely,
given $y$ in the right-hand side and $\epsilon>0$, choose a finite
$S\subset\Lambda$ with $n$ elements such that
$\sum_{\lambda\notin S}\lVert y_\lambda\rVert<\epsilon/2$ and, for $\lambda\in S$,
$z_\lambda\in\operatorname{span}(AX_\lambda)$ with
$\lVert y_\lambda-z_\lambda\rVert<\epsilon/(2n)$; the vector with these
coordinates on $S$ and $0$ elsewhere lies in $\operatorname{span}(AX)$ and is
within $\epsilon$ of $y$. Next, if $M_\lambda\subset X_\lambda$ are closed
subspaces and $M=\bigoplus_1 \{M_\lambda \}$, then
$\dist(x,M)=\sum_\lambda\dist(x_\lambda,M_\lambda)$: the inequality $\geq$ is
coordinatewise, and for $\leq$ one picks $m_\lambda\in M_\lambda$ with
$\lVert x_\lambda-m_\lambda\rVert\leq\dist(x_\lambda,M_\lambda)+\epsilon_\lambda$
for a summable family ${(\epsilon_\lambda)}$, the resulting $m$ being summable
because $\lVert m_\lambda\rVert\leq 2\lVert x_\lambda\rVert+\epsilon_\lambda$.
Together the two give the formula for $\sep_X$.

$(i)$ For $x\in X_{ann}$ each $x_\lambda$ lies in ${(X_\lambda)}_{ann}$, so
$\sep_{X_\lambda}(x_\lambda)\leq\sepd(X_\lambda)\lVert x_\lambda\rVert$ and
$(iv)$ gives
$\sep_X(x)\leq \bigl(\sup_\mu\sepd(X_\mu) \bigr)\lVert x\rVert$.

$(ii)$ For $p=\infty$ the inclusion
$X_{ess}\subset\bigoplus_\infty \{{(X_\lambda)}_{ess} \}$ holds by the argument
of $(iv)$, but the reverse can fail, since $\operatorname{cl}_X(AX)$ need not
reach every bounded family of essential vectors; and distances to a
$\bigoplus_\infty$-sum of closed subspaces dominate the supremum of the
coordinate distances. Distances to a larger set being smaller, this already gives
$\sep_X(x)\geq\sup_\lambda\sep_{X_\lambda}(x_\lambda)$. When $\Lambda$ is finite
there is no tail to discard, so the argument of $(iv)$ gives
$X_{ess}=\bigoplus_\infty \{{(X_\lambda)}_{ess} \}$, while a near-optimal choice
in each of the finitely many coordinates gives
$\dist(x,M)=\max_\lambda\dist(x_\lambda,M_\lambda)$; hence
$\sep_X(x)=\max_\lambda\sep_{X_\lambda}(x_\lambda)$, which for $x\in X_{ann}$ is
the equality in $(ii)$.

$(iii)$ Let $\Lambda$ be finite and ${(x^\nu)}$ an annihilator net in $B_X$. Pass
to a subnet along which $\sep_X(x^\nu)$ converges to its limit superior and
$\lVert x^\nu_\lambda\rVert\to t_\lambda$ for each of the finitely many
$\lambda$; then $\sum_\lambda t_\lambda\leq 1$ for $p=1$ and $t_\lambda\leq 1$
for $p=\infty$. If $t_\lambda=0$ then
$\sep_{X_\lambda}(x^\nu_\lambda)\leq\lVert x^\nu_\lambda\rVert\to0$. If
$t_\lambda>0$ then eventually $\lVert x^\nu_\lambda\rVert\geq t_\lambda/2>0$ and
the normalised vectors $x^\nu_\lambda/\lVert x^\nu_\lambda\rVert$ form an
annihilator net in $B_{X_\lambda}$, whence
$\limsup_\nu\sep_{X_\lambda}(x^\nu_\lambda)\leq t_\lambda\asepd(X_\lambda)$. Now
$(iv)$ for $p=1$ and $(ii)$ for $p=\infty$, together with the fact that the limit
superior of a finite sum, respectively of a finite maximum, is at most the sum,
respectively the maximum, of the limits superior, give
$\limsup_\nu\sep_X(x^\nu)\leq\sup_\lambda\asepd(X_\lambda)$. Finiteness enters
exactly here: it is what forces the masses $t_\lambda$ to account for the whole
net.
\end{proof}

\subsection{The separation of the unitization}\label{SubSectionSeparationOfUnitization}

\emph{Nothing of [B] is interleaved here.}

Evaluating the separation defect on the free module of rank one, $A_+$, and on
its dual produces two constants of the algebra. They differ exactly as an
identity differs from an approximate one, and the dual is the more generous of
the two for a structural reason: the dual of $A_+$ always has an annihilator
direction, namely the augmentation
$\varepsilon:A_+\to\mathbb{C},\,a\oplus_1\lambda\mapsto\lambda$, whereas $A_+$
itself has one only when $A$ has a right identity.

\begin{lemma}\label{AnnihPartOfUnitization}
Regard $A_+$ as a left $A$-module and $A_+^*$ as the dual right $A$-module. Then
${(A_+)}_{ess}=A$, so that $\sep_{A_+}(v)=|\varepsilon(v)|$ for every
$v\in A_+$; and a vector $y=b\oplus_1\lambda$ lies in ${(A_+)}_{ann}$ if and
only if either $\lambda=0$ and $Ab= \{0 \}$, or $\lambda\neq 0$ and
$u:=-b/\lambda$ is a right identity of $A$, in which case
$y=\lambda(e_{A_+}-u)$. Consequently
\[
\sepd(A_+)=\frac{1}{1+\cri(A)} ,
\]
which is positive if and only if $A$ has a right identity. Moreover
${(A_+^*)}_{ann}=A^\perp=\mathbb{C}\varepsilon$ with $\lVert\varepsilon\rVert=1$,
so that $\sepd(A_+^*)=\sep_{A_+^*}(\varepsilon)$: a single distance in a single
direction.
\end{lemma}

\begin{proof}
In $A_+=A\bigoplus_1\mathbb{C}$ the left action of $A$ is
$a\cdot(b\oplus_1\lambda)=(ab+\lambda a)\oplus_1 0$, so
$A\cdot A_+\subset A$; conversely $a=a\cdot e_{A_+}$, so
$\operatorname{span}(A\cdot A_+)=A$, which is closed. Hence
${(A_+)}_{ess}=A$ and, for $v=b\oplus_1\lambda$,
$\dist(v,A)=\inf_{c\in A}(\lVert b-c\rVert+|\lambda|)=|\varepsilon(v)|$.

The same formula shows that $y=b\oplus_1\lambda$ is annihilated by $A$ if and
only if $ab+\lambda a=0$ for every $a\in A$. If $\lambda=0$ this says
$Ab= \{0 \}$, and then $y\in A={(A_+)}_{ess}$, so $\sep_{A_+}(y)=0$. If
$\lambda\neq 0$, put $u:=-b/\lambda\in A$; the condition becomes $au=a$ for every
$a$, so $u$ is a right identity and $y=\lambda(e_{A_+}-u)$; conversely every
vector of that form is annihilated. For such a $y$,
$\sep_{A_+}(y)=|\lambda|$ while
$\lVert y\rVert=|\lambda|(1+\lVert u\rVert)$, so the ratio
$\sep_{A_+}(y)/\lVert y\rVert$ is $1/(1+\lVert u\rVert)$, the annihilator vectors
with $\lambda=0$ contributing $0$. Hence $\sepd(A_+)$ is the supremum of
$1/(1+\lVert u\rVert)$ over the right identities $u$ of $A$, which is
$1/(1+\cri(A))$, the supremum over the empty set being $0$ and matching the
convention $\cri(A)=+\infty$.

Finally, ${(A_+^*)}_{ann}={\bigl({(A_+)}_{ess}\bigr)}^\perp=A^\perp$; a
functional on $A\bigoplus_1\mathbb{C}$ vanishing on $A$ is determined by its
value at $e_{A_+}$, so $A^\perp=\mathbb{C}\varepsilon$, and
$\lVert\varepsilon\rVert=1$. The defect of $A_+^*$ is therefore the separation of
the single unit vector $\varepsilon$.
\end{proof}

The dual constant is computed by approximate identities instead, and there the
infimum is attained. No hypothesis on $A$ is needed: all three clauses hold for
the zero algebra, where $A_+=\mathbb{C}$ has ${(A_+)}_{ess}= \{0 \}$ and
${(A_+^*)}_{ess}= \{0 \}$, so that $\sepd(A_+^*)=\lVert\varepsilon\rVert=1$,
matching $1/(1+\crai( \{0 \}))=1$, and the zero net attains the bound.

\begin{theorem}\label{SepOfDualUnitization}
\begin{enumerate}[label = (\roman*)]
\item $\sep_{X^*}(f)\geq\lVert f\rVert/(1+\crai(A))$ for every left $A$-module
$X$ and every $f\in{(X_{ess})}^\perp$; equivalently
$\sepd(X^*)\geq\sepd(A_+^*)$ whenever $X$ is not essential;

\item $\sepd(A_+^*)=\dfrac{1}{1+\crai(A)}$;

\item $\crai(A)$, when finite, is attained by a right approximate identity.
\end{enumerate}
\end{theorem}

\begin{proof}
$(i)$ The dual $X^*$ is a right $A$-module with ${(X^*)}_{ann}={(X_{ess})}^\perp$,
so lemma~\ref{SepLowerBound} in its right-hand form, applied to $X^*$, gives the
displayed bound; if $A$ has no right approximate identity then $\crai(A)=+\infty$
and the bound is vacuous. For the reformulation, $\sepd(X^*)$ is the supremum of
$\sep_{X^*}(f)/\lVert f\rVert$ over $f\in{(X^*)}_{ann}\setminus \{0 \}$, and that
set is non-empty exactly when $X_{ess}\neq X$, by the Hahn-Banach theorem; the
bound then reads $\sepd(X^*)\geq 1/(1+\crai(A))$, which is $\sepd(A_+^*)$ by
$(ii)$, proved below and independently of $(i)$.

$(ii)$ By lemma~\ref{AnnihPartOfUnitization} we have
$\sepd(A_+^*)=\sep_{A_+^*}(\varepsilon)$, so
proposition~\ref{SepOfDualViaAnnihNets} at $X=A_+$ and $f=\varepsilon$ reads
\[
\sepd(A_+^*)=\sup \Bigl\{
  \textstyle\limsup_\nu|\lambda_\nu|:
  \lVert b_\nu\rVert+|\lambda_\nu|\leq 1,\
  \lim_\nu\lVert ab_\nu+\lambda_\nu a\rVert=0\ \text{for all }a\in A
 \Bigr\},
\]
the supremum being over nets $v_\nu=b_\nu\oplus_1\lambda_\nu$, since
$\varepsilon(v_\nu)=\lambda_\nu$ and
$a\cdot v_\nu=(ab_\nu+\lambda_\nu a)\oplus_1 0$. Call such a net admissible.

For the lower bound, suppose $A$ has a right approximate identity and let
$\epsilon>0$. Choose one, ${(e_\nu)}_{\nu\in N}$, with
$\sup_\nu\lVert e_\nu\rVert\leq\crai(A)+\epsilon$, and put
$v_\nu:=(e_{A_+}-e_\nu)/(1+\lVert e_\nu\rVert)$. These are unit vectors, hence
admissible; $a\cdot v_\nu=(a-ae_\nu)/(1+\lVert e_\nu\rVert)\to 0$; and
$\lambda_\nu=1/(1+\lVert e_\nu\rVert)\geq 1/(1+\crai(A)+\epsilon)$. Let
$\epsilon\to 0$.

For the upper bound, let ${(v_\nu)}$ be admissible with
$\limsup_\nu|\lambda_\nu|=d>0$; pass to a subnet along which
$|\lambda_\nu|\to d$, so that eventually $\lambda_\nu\neq 0$, and set
$u_\nu:=-b_\nu/\lambda_\nu\in A$. Then
\[
\lVert a-au_\nu\rVert=\frac{\lVert a\cdot v_\nu\rVert}{|\lambda_\nu|}
\longrightarrow 0\quad(a\in A),
\qquad
\lVert u_\nu\rVert=\frac{\lVert b_\nu\rVert}{|\lambda_\nu|}
\leq\frac{1-|\lambda_\nu|}{|\lambda_\nu|}\longrightarrow\frac{1-d}{d} ,
\]
so ${(u_\nu)}$ is a right approximate identity and, for every $\epsilon>0$, its
tail on which $\lVert u_\nu\rVert\leq(1-d)/d+\epsilon$ is one as well; hence
$\crai(A)\leq(1-d)/d$, that is $d\leq 1/(1+\crai(A))$. The same computation shows
that $d>0$ is impossible when $A$ has no right approximate identity, so
$\sepd(A_+^*)=0$ there, again the value of $1/(1+\crai(A))$.

$(iii)$ Assume $\crai(A)<+\infty$ and put $\delta:=\sepd(A_+^*)>0$. With
$E={(A_+^*)}_{ess}$, proposition~\ref{SepOfDualViaAnnihNets} says that the
distance $\delta=\dist(\varepsilon,E)$ is attained at some $q\in B_{A_+^{**}}$
with $q|_E=0$; multiplying $q$ by a unimodular scalar we may assume
$q(\varepsilon)=\delta$. Since $A_+=A\bigoplus_1\mathbb{C}$ we have
$A_+^*\isom{\mathbf{Ban}_1}A^*\bigoplus_\infty\mathbb{C}$ and
$A_+^{**}\isom{\mathbf{Ban}_1}A^{**}\bigoplus_1\mathbb{C}$, where $g\in A_+^*$ is
written $g=(h,c)$ with $h=g|_A$ and $c=g(e_{A_+})$, so that
$g=h+c\varepsilon$ and $\langle(m,t),(h,c)\rangle=m(h)+tc$. In these coordinates
$q=(m,\delta)$ with $\lVert m\rVert+\delta\leq 1$, that is
$\lVert m\rVert\leq 1-\delta$. The right action reads
\[
(g\cdot a)(b\oplus_1\lambda)=g \bigl((ab+\lambda a)\oplus_1 0 \bigr)
=(h\cdot a)(b)+\lambda\,h(a),
\]
that is $g\cdot a=\bigl(h\cdot a,\ h(a)\bigr)$, which does not involve $c$; since
$E$ is the closed span of the $g\cdot a$, the condition $q|_E=0$ says
$m(h\cdot a)+\delta h(a)=0$ for all $h\in A^*$ and $a\in A$. Setting
$p:=-m/\delta\in A^{**}$ we get
\[
p(h\cdot a)=h(a)\quad(h\in A^*,\ a\in A),
\qquad
\lVert p\rVert=\frac{\lVert m\rVert}{\delta}\leq\frac{1-\delta}{\delta}=\crai(A).
\]
By Goldstine's theorem applied to the ball of radius $\lVert p\rVert$ there is a
net ${(x_\mu)}$ in $A$ with $\lVert x_\mu\rVert\leq\lVert p\rVert$ and
$\iota_A(x_\mu)\to p$ weak${}^*$; then
$h(ax_\mu)=(h\cdot a)(x_\mu)\to p(h\cdot a)=h(a)$, so $ax_\mu\to a$ weakly for
every $a\in A$. A Mazur argument on the convex set
$ \{{(ax-a)}_{a\in F}:x\in A,\ \lVert x\rVert\leq\lVert p\rVert \}$, followed by
directing the pairs $(F,\epsilon)$ as in the proof of
proposition~\ref{SepOfDualViaAnnihNets}, produces a right approximate identity
bounded by $\lVert p\rVert\leq\crai(A)$; since $\crai(A)$ is the infimum of such
bounds, this one attains it.
\end{proof}

\begin{lemma}\label{AugmCharDistEstimate}
Let $A\neq \{0 \}$, let $a\in A$ with $\lVert a\rVert=1$, let
$\varphi\in A_+^*$ satisfy $\lVert\varphi\rVert=1$ and $\varphi(a)=1$, and let
$0<\gamma\leq 1$. Then
$\lVert\varepsilon-\gamma\,\varphi\cdot a\rVert\leq\max \{\gamma,1-\gamma \}$. In
particular $\sepd(A_+^*)\leq\tfrac{1}{2}$ for every $A\neq \{0 \}$.
\end{lemma}

\begin{proof}
A vector $a$ of norm one exists because $A\neq \{0 \}$, and such a $\varphi$
exists because on the line $\mathbb{C}a$ the functional $ta\mapsto t$ has norm
one and the Hahn-Banach theorem extends it to $A_+$. For
$b\oplus_1\lambda\in A_+$ one has
$a\cdot(b\oplus_1\lambda)=(ab+\lambda a)\oplus_1 0$, so, using $\varphi(a)=1$,
\[
(\varepsilon-\gamma\,\varphi\cdot a)(b\oplus_1\lambda)
=\lambda-\gamma\varphi(ab)-\gamma\lambda
=\lambda(1-\gamma)-\gamma\varphi(ab),
\]
whose modulus is at most
$(1-\gamma)|\lambda|+\gamma\lVert b\rVert
\leq\max \{\gamma,1-\gamma \} \bigl(|\lambda|+\lVert b\rVert \bigr)$. Since
$\varphi\cdot a\in{(A_+^*)}_{ess}$, lemma~\ref{AnnihPartOfUnitization} gives
$\sepd(A_+^*)=\sep_{A_+^*}(\varepsilon)
\leq\lVert\varepsilon-\gamma\,\varphi\cdot a\rVert
\leq\max \{\gamma,1-\gamma \}$, which at $\gamma=\tfrac{1}{2}$ is
$\tfrac{1}{2}$.
\end{proof}

\begin{corollary}\label{SepOfDualUnitizationValues}
Let $A\neq \{0 \}$. Then $\sepd(A_+^*)=\tfrac{1}{2}$ if and only if $A$ has a
contractive right approximate identity; $\sepd(A_+^*)=0$ if and only if $A$ has
no right approximate identity, if and only if
$\varepsilon\in{(A_+^*)}_{ess}$; and $\sepd(A_+^*)=1/(1+\lVert e_A\rVert)$ for
unital $A$.
\end{corollary}

\begin{proof}
By theorem~\ref{SepOfDualUnitization}$(ii)$ all three clauses are statements
about $\crai(A)$, which for $A\neq \{0 \}$ satisfies $\crai(A)\geq 1$ by
subsection~\ref{SubSectionBanachAlgebrasAndTheirModules} --- the same ceiling
$\sepd(A_+^*)\leq\tfrac{1}{2}$ that lemma~\ref{AugmCharDistEstimate} gives
independently. So $\sepd(A_+^*)=\tfrac{1}{2}$ if and only if $\crai(A)=1$, which
holds if $A$ has a contractive right approximate identity and, conversely, forces
one to exist by theorem~\ref{SepOfDualUnitization}$(iii)$. Next,
$\sepd(A_+^*)=0$ if and only if $\crai(A)=+\infty$, that is, if and only if $A$
has no right approximate identity; and
$\sepd(A_+^*)=\dist \bigl(\varepsilon,{(A_+^*)}_{ess} \bigr)$ vanishes exactly
when $\varepsilon$ lies in that closed subspace. The unital clause is
$\crai(A)=\lVert e_A\rVert$, recorded in
subsection~\ref{SubSectionBanachAlgebrasAndTheirModules}.

The hypothesis $A\neq \{0 \}$ is not decoration, and this is the only place in
this section where it bites: for $A= \{0 \}$ the conventions give
$\crai(A)=0$ while $\sepd(A_+^*)=1$, so the first clause would fail, the zero net
being a contractive right approximate identity.
\end{proof}

The dual constant runs over the whole of $(0,\tfrac{1}{2}]$. Let
$A_t:=(\mathbb{C},t|\cdot|)$, the field $\mathbb{C}$ with the norm
$\lVert z\rVert=t|z|$; this is a Banach algebra exactly for $t\geq 1$, since
$\lVert zw\rVert=t|zw|$ while $\lVert z\rVert\lVert w\rVert=t^2|zw|$. Its
identity is $1$, of norm $t$, so the unital clause of
corollary~\ref{SepOfDualUnitizationValues} gives
$\sepd \bigl({(A_t)}_+^* \bigr)=1/(1+t)$, and as $t$ runs over $[1,+\infty)$ this
runs over $(0,\tfrac{1}{2}]$. The identity map $A_t\to A_s$ is a topological
isomorphism of Banach algebras for all $s,t\geq 1$, so $\sepd(A_+^*)$ is an
isometric and not a topological invariant of $A$.

It depends on the side as well. Let
$R= \bigl\{ \bigl(\begin{smallmatrix}a&b\\0&0\end{smallmatrix} \bigr) \bigr\}
\subset M_2(\mathbb{C})$, a subalgebra because
$\bigl(\begin{smallmatrix}a&b\\0&0\end{smallmatrix} \bigr)
\bigl(\begin{smallmatrix}a'&b'\\0&0\end{smallmatrix} \bigr)
=\bigl(\begin{smallmatrix}aa'&ab'\\0&0\end{smallmatrix} \bigr)$. The matrix unit
$E_{11}$ lies in $R$, has norm one and satisfies $E_{11}x=x$ for every $x\in R$,
so $\cli(R)=\clai(R)=1$ and the mirror constant, computed from left approximate
identities, is $\tfrac{1}{2}$. On the other side, $E_{12}\in R$ satisfies
$E_{12}x=0$ for every $x\in R$, so no net ${(e_\nu)}$ in $R$ can have
$E_{12}e_\nu\to E_{12}\neq 0$: the algebra $R$ has no right approximate identity,
$\crai(R)=+\infty$, and $\sepd \bigl(R_+^* \bigr)=0$.

One reading of proposition~\ref{MetTopFlatCycModCharac} of [B] is worth
recording here, since it is the first place in [B] where the dual constant
occurs unnamed. For a closed left ideal $I$ of $A_+$ the quantity
$\sup_\nu\lVert e_{A_+}-e_\nu\rVert$ appearing there equals
$1+\sup_\nu\lVert e_\nu\rVert$, because $e_\nu\in I\subset A$; its infimum over
the right bounded approximate identities of $I$ is therefore
$1+\crai(I)=1/\sepd \bigl(I_+^* \bigr)$, and that proposition is the computation
of this constant along the coset of the identity.
Proposition~\ref{SepOfDualViaAnnihNets} and
theorem~\ref{SepOfDualUnitization} give it directly.

Both constants are constants of $A_+$ alone: free modules of higher rank add
nothing to either. For $\sepd$ this is proposition~\ref{SepOfSums}$(i)$, a free
module being an $\ell_1$-sum of copies of $A_+$; the asymptotic defect needs an
argument.

\begin{proposition}\label{AsympSepOfFreeMod}
$\asepd \bigl(A_+\projtens\ell_1(\Lambda) \bigr)=\sepd(A_+^*)$ for every
$\Lambda\neq\varnothing$, a value at most $\tfrac{1}{2}$ when $A\neq \{0 \}$:
free modules have the minimal asymptotic separation defect.
\end{proposition}

\begin{proof}
First, $\asepd(A_+)=\sepd(A_+^*)$. Indeed
$\sep_{A_+}(v)=|\varepsilon(v)|$ by lemma~\ref{AnnihPartOfUnitization}, so an
annihilator net ${(v_\nu)}$ in $B_{A_+}$ has
$\asep_{A_+} \bigl({(v_\nu)} \bigr)=\limsup_\nu|\varepsilon(v_\nu)|$; taking the
supremum over such nets, definition~\ref{AsympSepOfMod} and
proposition~\ref{SepOfDualViaAnnihNets} at $(A_+,\varepsilon)$ give
$\asepd(A_+)=\sep_{A_+^*}(\varepsilon)=\sepd(A_+^*)$.

Write $P=A_+\projtens\ell_1(\Lambda)$. The canonical map carrying
$u=\sum_\lambda u_\lambda\otimes\delta_\lambda$ to ${(u_\lambda)}$ is an
isometric isomorphism of left $A$-modules onto
$\bigoplus_1 \{A_+:\lambda\in\Lambda \}$, so by
proposition~\ref{SepOfSums}$(iv)$ and lemma~\ref{AnnihPartOfUnitization} the
essential part of $P$ is $A\projtens\ell_1(\Lambda)$ and, writing
$u_\lambda=b_\lambda\oplus_1\beta_\lambda$,
\[
\sep_P(u)=\sum_\lambda\sep_{A_+}(u_\lambda)=\sum_\lambda|\beta_\lambda| .
\]
A coordinate copy of $A_+$ is a $1$-retract of $P$, so
lemma~\ref{SepUnderRetracts} gives $\asepd(P)\geq\asepd(A_+)$.

For the upper bound let ${(u^\nu)}_{\nu\in N}$ be an annihilator net in $B_P$.
For each $\nu$ and $\lambda$ choose a unimodular $c^\nu_\lambda$ with
$c^\nu_\lambda\beta^\nu_\lambda=|\beta^\nu_\lambda|$ and put
$w_\nu:=\sum_\lambda c^\nu_\lambda u^\nu_\lambda\in A_+$, a series converging
absolutely because $\sum_\lambda\lVert u^\nu_\lambda\rVert=\lVert u^\nu\rVert
\leq 1$. Then $\lVert w_\nu\rVert\leq 1$ and
$\lVert a\cdot w_\nu\rVert\leq\sum_\lambda\lVert a\cdot u^\nu_\lambda\rVert
=\lVert a\cdot u^\nu\rVert\to 0$, so ${(w_\nu)}$ is an annihilator net in
$B_{A_+}$; and, $\varepsilon$ being linear and continuous,
$\varepsilon(w_\nu)=\sum_\lambda|\beta^\nu_\lambda|=\sep_P(u^\nu)$. Hence
$\limsup_\nu\sep_P(u^\nu)=\limsup_\nu|\varepsilon(w_\nu)|$, which
proposition~\ref{SepOfDualViaAnnihNets} at $(A_+,\varepsilon)$ bounds by
$\sepd(A_+^*)$. Taking the supremum over annihilator nets in $B_P$ gives the
stated equality, and the bound $\tfrac{1}{2}$ for $A\neq \{0 \}$ is
lemma~\ref{AugmCharDistEstimate}.
\end{proof}

The dual of the free module of rank $\Lambda$ is the cofree module of rank
$\Lambda$, and its separation defect is the second constant, not the first. The
upper bound comes from the predual: that is the only use the chain of
proposition~\ref{TheChain} gets outside subsection~\ref{SubSectionSeparation}.

\begin{proposition}\label{SepOfCofreeMod}
For every left $A$-module $X$ that is not essential,
\[
\sepd(A_+^*)\ \leq\ \sepd(X^*)\ \leq\ \asepd(X) ,
\]
so $\sepd(X^*)=\sepd(A_+^*)$ as soon as $X$ attains the minimal asymptotic
separation defect. In particular, for every $\Lambda\neq\varnothing$ the cofree
right $A$-module satisfies
\[
\sepd \bigl(\mathcal{B}(A_+,\ell_\infty(\Lambda)) \bigr)=\sepd(A_+^*)
=\frac{1}{1+\crai(A)} .
\]
\end{proposition}

\begin{proof}
If $X$ is not essential then ${(X_{ess})}^\perp\neq \{0 \}$ by the Hahn-Banach
theorem, so theorem~\ref{SepOfDualUnitization}$(i)$ gives
$\sepd(X^*)\geq\sepd(A_+^*)$, while proposition~\ref{TheChain} gives
$\sepd(X^*)\leq\asepd(X)$; if $\asepd(X)=\sepd(A_+^*)$ the outer terms coincide.

Write $F_\Lambda=A_+\projtens\ell_1(\Lambda)$. Since
$\ell_\infty(\Lambda)\isom{\mathbf{Ban}_1}{\ell_1(\Lambda)}^*$, the law of
adjoint associativity gives isometric isomorphisms of right $A$-modules
\[
\mathcal{B} \bigl(A_+,\ell_\infty(\Lambda) \bigr)
\isom{\mathbf{mod}_1\text{-}A}
\mathcal{B} \bigl(A_+,\mathcal{B}(\ell_1(\Lambda),\mathbb{C}) \bigr)
\isom{\mathbf{mod}_1\text{-}A}
\mathcal{B} \bigl(A_+\projtens\ell_1(\Lambda),\mathbb{C} \bigr)=F_\Lambda^* .
\]
By proposition~\ref{SepOfSums}$(iv)$ and lemma~\ref{AnnihPartOfUnitization} the
essential part of $F_\Lambda$ is $A\projtens\ell_1(\Lambda)$, a proper closed
subspace since $A$ is proper in $A_+$ and $\Lambda\neq\varnothing$; so
$F_\Lambda$ is not essential, and
proposition~\ref{AsympSepOfFreeMod} gives $\asepd(F_\Lambda)=\sepd(A_+^*)$. The
two bounds therefore meet at $X=F_\Lambda$, and the value is
$1/(1+\crai(A))$ by theorem~\ref{SepOfDualUnitization}$(ii)$.
\end{proof}

So the free modules see the identities of $A$ and the cofree ones its
approximate identities:
$\sepd(F_\Lambda)=\sepd(A_+)=1/(1+\cri(A))$ by
proposition~\ref{SepOfSums}$(i)$ and lemma~\ref{AnnihPartOfUnitization}, against
$\sepd(F_\Lambda^*)=1/(1+\crai(A))$. Corollary~\ref{SepOfInjMod} will recover the
second value in one line and for all metrically injective modules at once.

\subsection{The annihilator part of a projective module}\label{SubSectionAnnihPartOfProjMod}

\emph{Interleaved from [B], in this order, before the material below:}
proposition~\ref{MetTopRelProjModCompIdealPartCompl},
corollary~\ref{MetTopRelProjModPartCompl},
proposition~\ref{MetTopProjModEssQuotIsProjBanSp}.

Those three say that the essential part of a projective module is complemented
and that the quotient by it is a projective Banach space. What follows measures
how far the annihilator part stands out of the essential one, and it does so
through a single lifting problem: the augmentation $\varepsilon:A_+\to\mathbb{C}$.
The retract characterisation~\ref{MetCTopProjModViaCanonicMorph} is not needed
anywhere in this subsection.

\begin{lemma}\label{ProjLiftAlongAugmentation}
Let $P$ be a $\langle$~metrically / $C$-topologically~$\rangle$ projective left
$A$-module. Then for every $x\in P$ there is an $A$-morphism $\psi:P\to A_+$ with
$\lVert\psi\rVert\leq\langle\,1\ /\ C\,\rangle$ and
$\varepsilon(\psi(x))=\sep_P(x)$.
\end{lemma}

\begin{proof}
Regard $\mathbb{C}$ as a left $A$-module with the zero action. By the
Hahn-Banach description of $\sep_P(x)=\dist(x,P_{ess})$ there is
$g\in{(P_{ess})}^\perp$ with $\lVert g\rVert\leq 1$ and $g(x)=\sep_P(x)$; take
$g=0$ if $x\in P_{ess}$. Then $g:P\to\mathbb{C}$ is an $A$-morphism, since
$a\cdot y\in P_{ess}$ gives $g(a\cdot y)=0=a\cdot g(y)$. For the same two reasons
$\varepsilon$ is an $A$-morphism,
$\varepsilon(a\cdot v)=0=a\cdot\varepsilon(v)$; it has norm one and carries
$B_{A_+}$ onto $B_{\mathbb{C}}$, since $\varepsilon(0\oplus_1\mu)=\mu$ with
$\lVert 0\oplus_1\mu\rVert=|\mu|$, so it is strictly coisometric and in
particular strictly $1$-topologically surjective. By
definition~\ref{MetCTopProjMod}, with $c=1$ in the topological bracket,
projectivity of $P$ lifts $g$ along $\varepsilon$: there is an $A$-morphism
$\psi:P\to A_+$ with $\varepsilon\psi=g$ and
$\lVert\psi\rVert\leq\langle\,\lVert g\rVert\ /\ C\lVert g\rVert\,\rangle
\leq\langle\,1\ /\ C\,\rangle$. Finally
$\varepsilon(\psi(x))=g(x)=\sep_P(x)$.
\end{proof}

\begin{theorem}\label{SepOfProjMod}
Let $P$ be a $\langle$~metrically / $C$-topologically~$\rangle$ projective left
$A$-module. Then
\[
\sep_P(x)\leq\langle\,\sepd(A_+)\ /\ C\,\sepd(A_+)\,\rangle\lVert x\rVert
\qquad(x\in P_{ann}),
\]
that is, $\sepd(P)\leq\langle\,\sepd(A_+)\ /\ C\sepd(A_+)\,\rangle$.
\end{theorem}

\begin{proof}
Fix $x\in P_{ann}$ and let $\psi$ be as in
lemma~\ref{ProjLiftAlongAugmentation}. The vector $\psi(x)$ lies in
${(A_+)}_{ann}$, because $a\cdot\psi(x)=\psi(a\cdot x)=0$ for every $a\in A$, so
lemma~\ref{AnnihPartOfUnitization} and the definition of $\sepd$ give
\[
\sep_P(x)=\varepsilon(\psi(x))=\sep_{A_+} \bigl(\psi(x) \bigr)
\leq\sepd(A_+)\lVert\psi(x)\rVert
\leq\langle\,\sepd(A_+)\ /\ C\,\sepd(A_+)\,\rangle\lVert x\rVert .
\]
Dividing by $\lVert x\rVert$ and taking the supremum over
$x\in P_{ann}\setminus \{0 \}$ gives the statement about $\sepd(P)$.
\end{proof}

One lifting problem suffices, and it is the one dual to the extension problem
that governs injectivity: projectivity meets the annihilator module
$\mathbb{C}\isom{A\text{-}\mathbf{mod}_1}A_+/A$ as a \emph{quotient} of $A_+$ and
lifts along the augmentation, while injectivity and flatness meet its dual
$\mathbb{C}\varepsilon=A^\perp$ as a \emph{subobject} of $A_+^*$ and extend along
the inclusion or tensor with it --- see
propositions~\ref{SepOfAnnihVectorInInjMod} and~\ref{SepViaAnnihVector}$(ii)$.
That the projective constant is the smaller of the two ---
proposition~\ref{SepVsSepDual} at $X=A_+$ --- is nothing but the difference
between splitting the augmentation sequence exactly and splitting it
approximately.

\begin{corollary}\label{ProjAnnihMod}
Let $P\neq \{0 \}$ be an annihilator $A$-module which is
$\langle$~metrically / $C$-topologically~$\rangle$ projective. Then
$\langle$~$A= \{0 \}$ / $\cri(A)\leq C-1$~$\rangle$.
\end{corollary}

\begin{proof}
Since $P$ is a non-zero annihilator module, $P_{ess}= \{0 \}$ and $P_{ann}=P$, so
$\sep_P(x)=\lVert x\rVert$ for every $x\in P$ and $\sepd(P)=1$.
Theorem~\ref{SepOfProjMod} therefore gives
$\sepd(A_+)\geq\langle\,1\ /\ 1/C\,\rangle$, and
lemma~\ref{AnnihPartOfUnitization} turns this into a statement about $\cri(A)$
through $\sepd(A_+)=1/(1+\cri(A))$. In the topological bracket it reads
$\cri(A)\leq C-1$. In the metric bracket it reads $\cri(A)\leq 0$; but
$\cri(A)\geq 1$ for $A\neq \{0 \}$ by
subsection~\ref{SubSectionBanachAlgebrasAndTheirModules}, so only $A= \{0 \}$ is
compatible with it.
\end{proof}

Corollary~\ref{ProjAnnihMod} is stated as a necessary condition only. Its
converse is the implication $(ii)\Rightarrow(i)$ of
proposition~\ref{MetTopProjOfAnnihModCharac}, into which the corollary is
dissolved in subsection~\ref{SubSectionHomoligicallyTrivialAnnihilatorModules};
it is recorded here because it is what
theorem~\ref{SepOfProjMod} contributes, and it is not cited anywhere else in
this subsection.

\begin{corollary}\label{ProjOverNonUnital}
Let $A$ have no right identity. Then $\sepd(A_+)=0$, so in every topologically
projective $A$-module $P$ one has $P_{ann}\subset P_{ess}$, and no non-zero
annihilator module is projective, metrically or topologically.
\end{corollary}

\begin{proof}
If $A$ has no right identity then $\cri(A)=+\infty$, so $\sepd(A_+)=0$ by
lemma~\ref{AnnihPartOfUnitization}. Let $P$ be $C$-topologically projective for
some $C$ and let $x\in P_{ann}$; theorem~\ref{SepOfProjMod} gives
$\sep_P(x)\leq C\,\sepd(A_+)\lVert x\rVert=0$, and $P_{ess}$ is closed, so
$x\in P_{ess}$.

For the last assertion let $P\neq \{0 \}$ be an annihilator module. Then
$P_{ess}= \{0 \}$ and $P_{ann}=P$, so $\sepd(P)=1$, and
theorem~\ref{SepOfProjMod} forces
$\sepd(A_+)\geq\langle\,1\ /\ 1/C\,\rangle>0$ in either bracket, contradicting
$\sepd(A_+)=0$.
\end{proof}

On the asymptotic layer the answer changes, and it is the dual constant that
appears.

\begin{proposition}\label{AsympSepOfProjMod}
Let $A\neq \{0 \}$ and let $P$ be a
$\langle$~metrically / $C$-topologically~$\rangle$ projective left $A$-module. If
$P$ is essential, then $\asepd(P)=0$; otherwise
\[
\langle\,\asepd(P)=\sepd(A_+^*)\ /\
\sepd(A_+^*)\leq\asepd(P)\leq C\,\sepd(A_+^*)\,\rangle .
\]
\end{proposition}

\begin{proof}
If $P$ is essential then $\sep_P$ vanishes identically and $\asepd(P)=0$.

Let $P$ be non-essential, let ${(x_\nu)}_{\nu\in N}$ be an annihilator net in
$B_P$ and write $M:=\langle\,1\ /\ C\,\rangle$, a number at least $1$. For each
$\nu$ let $\psi_\nu$ be the morphism supplied by
lemma~\ref{ProjLiftAlongAugmentation} at $x_\nu$ and put
$v_\nu:=\psi_\nu(x_\nu)$. Then $\lVert v_\nu\rVert\leq M$ and
$\lVert a\cdot v_\nu\rVert=\lVert\psi_\nu(a\cdot x_\nu)\rVert
\leq M\lVert a\cdot x_\nu\rVert\to 0$, so ${(v_\nu/M)}$ is an annihilator net in
$B_{A_+}$, while $\varepsilon(v_\nu)=\sep_P(x_\nu)$.
Proposition~\ref{SepOfDualViaAnnihNets} at $(A_+,\varepsilon)$ therefore gives
$\limsup_\nu\sep_P(x_\nu)/M\leq\sep_{A_+^*}(\varepsilon)=\sepd(A_+^*)$, and
taking the supremum over annihilator nets in $B_P$ gives
$\asepd(P)\leq\langle\,\sepd(A_+^*)\ /\ C\,\sepd(A_+^*)\,\rangle$.

For the lower bound, $P$ is not essential, so proposition~\ref{TheChain} and
theorem~\ref{SepOfDualUnitization}$(i)$ give
$\asepd(P)\geq\sepd(P^*)\geq\sepd(A_+^*)$. In the metric bracket the two bounds
coincide.
\end{proof}

Comparing this with theorem~\ref{SepOfProjMod} shows that the two layers must be
kept apart. Over $c_0(\Lambda)$ a non-essential projective module has
$\sepd(P)=0$ and $\asepd(P)=\tfrac{1}{2}$: every annihilator vector lies inside
$P_{ess}$, and the whole defect is carried by annihilator nets in $B_P$ that stay
at positive distance from $P_{ann}$.

\subsection{The annihilator part of an injective module}\label{SubSectionAnnihPartOfInjMod}

\emph{Interleaved from [B], in this order, before the material below:}
proposition~\ref{MetTopRelInjModCompIdealAnnihCompl},
corollary~\ref{MetTopRelInjModAnnihCompl},
proposition~\ref{MetTopInjModAnnihPartIsInjBanSp}.

Injectivity and flatness are governed by the dual constant instead, and by the
same mechanism: an annihilator direction of a dual module is transported into the
module under test. For injectivity the transport is an extension, for flatness a
tensor factor. Lemma~\ref{SepLowerBound} is the pointwise content of the
averaging projection built in the proof of
proposition~\ref{MetTopRelInjModCompIdealAnnihCompl}: the annihilator part is not
merely complemented, each of its vectors stands at a definite distance from the
essential part.

\begin{proposition}\label{SepOfAnnihVectorInInjMod}
Let $J$ be a $\langle$~metrically / $C$-topologically~$\rangle$ injective right
$A$-module, let $X$ be a left $A$-module and let $f\in{(X_{ess})}^\perp$ with
$\lVert f\rVert=1$. Then
\[
\sep_J(f_0)\leq\langle\,1\ /\ C\,\rangle\,\sep_{X^*}(f)\,\lVert f_0\rVert
\qquad(f_0\in J_{ann});
\]
in particular
$\sep_J(f_0)\leq\langle\,\sepd(A_+^*)\ /\ C\sepd(A_+^*)\,\rangle\lVert f_0\rVert$.
\end{proposition}

\begin{proof}
Since $f\in{(X_{ess})}^\perp={(X^*)}_{ann}$, the line
$\mathbb{C}f\subset X^*$ is a closed submodule on which the action is zero, and
the inclusion $\xi:\mathbb{C}f\to X^*$ is an isometric $A$-morphism, hence
$1$-topologically injective. The map
$\varphi:\mathbb{C}f\to J,\,\lambda f\mapsto\lambda f_0$ is well defined and
linear with $\lVert\varphi\rVert=\lVert f_0\rVert$, and it is an $A$-morphism:
$\varphi \bigl((\lambda f)\cdot a \bigr)=\varphi(0)=0$ while
$\varphi(\lambda f)\cdot a=\lambda(f_0\cdot a)=0$ because $f_0\in J_{ann}$. By
definition~\ref{MetCTopInjMod}, with $c=1$, injectivity of $J$ extends $\varphi$
along $\xi$: there is an $A$-morphism $\psi:X^*\to J$ with $\psi(f)=f_0$ and
$\lVert\psi\rVert\leq\langle\,\lVert f_0\rVert\ /\ C\lVert f_0\rVert\,\rangle$.

A morphism carries essential parts into essential parts, since
$\psi(g\cdot a)=\psi(g)\cdot a\in\operatorname{span}(JA)$ and $J_{ess}$ is
closed; so $\psi \bigl({(X^*)}_{ess} \bigr)\subset J_{ess}$ and, for every
$g\in{(X^*)}_{ess}$,
\[
\sep_J(f_0)\leq\lVert f_0-\psi(g)\rVert=\lVert\psi(f)-\psi(g)\rVert
\leq\lVert\psi\rVert\lVert f-g\rVert .
\]
Taking the infimum over $g$ turns the right-hand side into
$\lVert\psi\rVert\sep_{X^*}(f)$, which is the displayed bound. For the last
assertion take $(X,f)=(A_+,\varepsilon)$, admissible by
lemma~\ref{AnnihPartOfUnitization}, which also gives
$\sep_{A_+^*}(\varepsilon)=\sepd(A_+^*)$.
\end{proof}

\begin{corollary}\label{ExactSepOfAnnihVectorInInjMod}
Let $A$ have a right approximate identity. Then in every metrically injective
right $A$-module $J$,
\[
\sep_J(f_0)=\sepd(A_+^*)\lVert f_0\rVert=\frac{\lVert f_0\rVert}{1+\crai(A)}
\qquad(f_0\in J_{ann}),
\]
and $\sepd(A_+^*)\lVert f_0\rVert\leq\sep_J(f_0)
\leq C\sepd(A_+^*)\lVert f_0\rVert$ in every $C$-topologically injective one. The
value is attained on $A_+^*$ at $f_0=\varepsilon$. If $A$ has no right
approximate identity, then $J_{ann}\subset J_{ess}$ in every topologically
injective right $A$-module.
\end{corollary}

\begin{proof}
The upper bounds are proposition~\ref{SepOfAnnihVectorInInjMod} with
$(X,f)=(A_+,\varepsilon)$. For the lower bound, let $\epsilon>0$ and let
${(e_\nu)}$ be a right approximate identity bounded by $K=\crai(A)+\epsilon$;
lemma~\ref{SepLowerBound} in its right-hand form, applied to $J$, gives
$\sep_J(f_0)\geq\lVert f_0\rVert/(1+K)$, and letting $\epsilon\to 0$ and using
theorem~\ref{SepOfDualUnitization}$(ii)$ gives
$\sep_J(f_0)\geq\sepd(A_+^*)\lVert f_0\rVert$. In the metric bracket the two
bounds agree.

The value is attained: $A_+$ is free of rank one, hence metrically projective by
proposition~\ref{MetCTopFreeMod}, so $A_+^*$ is metrically injective by
proposition~\ref{DualMetTopProjIsMetrInj}, and its annihilator part is the line
$\mathbb{C}\varepsilon$ with $\lVert\varepsilon\rVert=1$, on which
lemma~\ref{AnnihPartOfUnitization} gives
$\sep_{A_+^*}(\varepsilon)=\sepd(A_+^*)$.

Finally, suppose $A$ has no right approximate identity; then $A\neq \{0 \}$,
since the zero net is one for the zero algebra, so
corollary~\ref{SepOfDualUnitizationValues} gives $\sepd(A_+^*)=0$ and the
topological upper bound reads $\sep_J(f_0)\leq 0$ for every $f_0\in J_{ann}$. As
$J_{ess}$ is closed, this says $J_{ann}\subset J_{ess}$.
\end{proof}

\begin{corollary}\label{InjAnnihModApprIdent}
Let $J\neq \{0 \}$ be an annihilator right $A$-module which is
$\langle$~metrically / $C$-topologically~$\rangle$ injective. Then
$\langle$~$A= \{0 \}$ / $A$ has a right $(C-1)$-bounded approximate
identity~$\rangle$.
\end{corollary}

\begin{proof}
Since $J$ is a non-zero annihilator module, $J_{ess}= \{0 \}$ and $J_{ann}=J$, so
$\sep_J(f_0)=\lVert f_0\rVert$ for every $f_0\in J$. Applying
proposition~\ref{SepOfAnnihVectorInInjMod} at some $f_0\neq 0$ with
$(X,f)=(A_+,\varepsilon)$ gives
$\sepd(A_+^*)\geq\langle\,1\ /\ 1/C\,\rangle$. In the metric bracket this says
$\sepd(A_+^*)\geq 1$, contradicting the ceiling $\tfrac{1}{2}$ of
lemma~\ref{AugmCharDistEstimate} unless the hypothesis of that lemma fails, that
is unless $A= \{0 \}$. In the topological bracket it says
$1/(1+\crai(A))\geq 1/C$, that is $\crai(A)\leq C-1$; by
theorem~\ref{SepOfDualUnitization}$(iii)$ this bound is attained, so $A$ has a
right approximate identity bounded by $C-1$.
\end{proof}

Corollary~\ref{ExactSepOfAnnihVectorInInjMod} is pointwise, and taking the
supremum over $B_{J_{ann}}$ turns it into a statement about the defect itself: on
metrically injective modules $\sepd$ is two-valued.

\begin{corollary}\label{SepOfInjMod}
Let $J$ be a $\langle$~metrically / $C$-topologically~$\rangle$ injective right
$A$-module. If $J$ is faithful, then $\sepd(J)=0$; otherwise
\[
\langle\,\sepd(J)=\sepd(A_+^*)\ /\
\sepd(A_+^*)\leq\sepd(J)\leq C\,\sepd(A_+^*)\,\rangle ,
\]
so that on metrically injective modules the separation defect takes only the two
values $0$ and $\sepd(A_+^*)$. Consequently, in the metric case:
\begin{enumerate}[label = (\roman*)]
\item $\sepd(P^*)=\sepd(A_+^*)$ for every left $A$-module $P$ that is metrically
projective and not essential --- the cofree modules of
proposition~\ref{SepOfCofreeMod} being the case
$P=A_+\projtens\ell_1(\Lambda)$;

\item $\sepd \bigl(\bigoplus_\infty \{J_\gamma:\gamma\in\Gamma \} \bigr)
=\sup_\gamma\sepd(J_\gamma)$ for every family of metrically injective right
$A$-modules, where proposition~\ref{SepOfSums}$(ii)$ gives only an inequality
once the family is infinite.
\end{enumerate}
\end{corollary}

\begin{proof}
Faithfulness means $J_{ann}= \{0 \}$, and then $\sepd(J)=\sep_J(0)=0$. Assume
$J_{ann}\neq \{0 \}$, so that $J_{ann}$ contains a vector of norm one. If $A$ has
a right approximate identity, corollary~\ref{ExactSepOfAnnihVectorInInjMod} gives
the pointwise equality, respectively the two pointwise bounds, and taking the
supremum over $f_0\in B_{J_{ann}}$ gives the displayed bracket, the supremum
being attained at a vector of norm one. If $A$ has no right approximate identity
then $A\neq \{0 \}$, so corollary~\ref{SepOfDualUnitizationValues} gives
$\sepd(A_+^*)=0$, while the last clause of
corollary~\ref{ExactSepOfAnnihVectorInInjMod} gives $J_{ann}\subset J_{ess}$ and
hence $\sepd(J)=0$; both brackets then read $0=0$.

$(i)$ Let $P$ be metrically projective and not essential. Its dual is metrically
injective by proposition~\ref{DualMetTopProjIsMetrInj}, and
${(P^*)}_{ann}={(P_{ess})}^\perp$ is non-zero precisely because
$P_{ess}\neq P$; so the main clause gives $\sepd(P^*)=\sepd(A_+^*)$. The cofree
modules are the case $P=A_+\projtens\ell_1(\Lambda)$, which is metrically
projective by proposition~\ref{MetCTopFreeMod} and whose dual is
$\mathcal{B}(A_+,\ell_\infty(\Lambda))$ by the identification made in the proof
of proposition~\ref{SepOfCofreeMod}.

$(ii)$ The sum is metrically injective by
proposition~\ref{MetTopInjModProd}, and each $J_\gamma$ is a $1$-retract of it,
so lemma~\ref{SepUnderRetracts} gives $\geq$. For the reverse, apply the main
clause to the sum itself: by proposition~\ref{SepOfSums} its annihilator part is
$\bigoplus_\infty \{{(J_\gamma)}_{ann}:\gamma\in\Gamma \}$, which is non-zero
exactly when some ${(J_\gamma)}_{ann}$ is. If none is, both sides are $0$. If
some is, the left-hand side equals $\sepd(A_+^*)$, and so does $\sepd(J_\gamma)$
for that $\gamma$.
\end{proof}

\subsection{The annihilator part of a flat module}\label{SubSectionAnnihPartOfFlatMod}

\emph{Interleaved from [B], in this order, before the material below:}
proposition~\ref{MetTopRelFlatModCompIdealPartCompl},
corollary~\ref{MetTopRelFlatModPartCompl},
proposition~\ref{MetTopFlatModEssQuotIsFlatBanSp}.

The engine here is not the separation defect but a statement about the module
tensor product; the defect appears in its corollaries.

\begin{proposition}\label{SepViaAnnihVector}
Let $F$ be a left $A$-module.
\begin{enumerate}[label = (\roman*)]
\item For every right $A$-module $Y$, every $y_0\in Y_{ann}$ and every $x\in F$,
\[
\lVert y_0\otimes x\rVert_{Y\projmodtens{A}F}\leq\lVert y_0\rVert\,\sep_F(x).
\]

\item If $F$ is $\langle$~metrically / $C$-topologically~$\rangle$ flat, then
\[
\bigl\langle\,
\lVert y_0\otimes x\rVert_{Y\projmodtens{A}F}=\lVert y_0\rVert\,\sep_F(x)
\ /\
\lVert y_0\rVert\,\sep_F(x)\leq C\lVert y_0\otimes x\rVert_{Y\projmodtens{A}F}
\,\bigr\rangle .
\]

\item Consequently, if $F$ is $\langle$~metrically / $C$-topologically~$\rangle$
flat, then for every left $A$-module $X$, every $f\in{(X_{ess})}^\perp$ with
$\lVert f\rVert=1$, every $x\in F$ and all finite families
$g_1,\dots,g_n\in X^*$ and $a_1,\dots,a_n\in A$,
\[
\sep_F(x)\leq\langle\,1\ /\ C\,\rangle\cdot \Bigl(
  \bigl\lVert f-\textstyle\sum_i g_i\cdot a_i\bigr\rVert\lVert x\rVert
  +\textstyle\sum_i\lVert g_i\rVert\lVert a_i\cdot x\rVert
 \Bigr).
\]
\end{enumerate}
\end{proposition}

\begin{proof}
$(i)$ By the law of adjoint associativity there is an isometric isomorphism
$\mathcal{B}(Y\projmodtens{A}F,\mathbb{C})
\isom{\mathbf{Ban}_1}\mathcal{B}_A(Y,F^*)$, under which a functional $\Phi$
corresponds to the morphism $T$ with
$\langle\Phi,y\otimes x\rangle=\langle Ty,x\rangle$. The norm of a vector being
the supremum of the moduli of the unit functionals at it,
\[
\lVert y_0\otimes x\rVert_{Y\projmodtens{A}F}
=\sup \{|\langle Ty_0,x\rangle|:T\in\mathcal{B}_A(Y,F^*),\,\lVert T\rVert\leq 1 \}.
\]
Fix such a $T$. It is a morphism into the right module $F^*$, so
$(Ty_0)\cdot a=T(y_0\cdot a)=0$ for every $a\in A$; that is,
$Ty_0\in{(F^*)}_{ann}={(F_{ess})}^\perp$ with
$\lVert Ty_0\rVert\leq\lVert y_0\rVert$. Hence
$|\langle Ty_0,x\rangle|=|\langle Ty_0,x-z\rangle|\leq\lVert y_0\rVert
\lVert x-z\rVert$ for every $z\in F_{ess}$, and the infimum over $z$ is
$\lVert y_0\rVert\sep_F(x)$.

$(ii)$ Both sides are positively homogeneous of degree one in $y_0$ and vanish at
$y_0=0$, so we may assume $\lVert y_0\rVert=1$. The line $Y_0:=\mathbb{C}y_0$ is
a closed submodule of $Y$ on which the action is zero, and the inclusion
$\xi:Y_0\to Y$ is an isometric $A$-morphism, hence also $1$-topologically
injective, so it belongs to the test class of definition~\ref{MetCTopFlatMod} in
both brackets with $c=1$; flatness of $F$ therefore makes
$\xi\projmodtens{A}1_F$ $\langle$~isometric / $C$-topologically
injective~$\rangle$. Now identify the source. The assignment
$y_0\otimes x\mapsto x$ extends to an isometric isomorphism
$Y_0\projtens F\to F$, since $Y_0$ is one-dimensional and spanned by a unit
vector; under it the subspace by which one divides to form
$Y_0\projmodtens{A}F$, namely
$\operatorname{cl}(\operatorname{span} \{
  y_0\cdot a\otimes x-y_0\otimes a\cdot x \})
=\operatorname{cl}(\operatorname{span} \{y_0\otimes a\cdot x \})$, corresponds to
$F_{ess}$. Passing to quotients gives an isometric isomorphism
$Y_0\projmodtens{A}F\isom{\mathbf{Ban}_1}F/F_{ess}$ carrying $y_0\otimes x$ to
the coset $x+F_{ess}$, so that
$\lVert y_0\otimes x\rVert_{Y_0\projmodtens{A}F}=\sep_F(x)$. Feeding this through
$\xi\projmodtens{A}1_F$ gives $(ii)$ in both brackets.

$(iii)$ Apply $(ii)$ with $(Y,y_0)=(X^*,f)$, legitimate because
$f\in{(X_{ess})}^\perp={(X^*)}_{ann}$ and $\lVert f\rVert=1$; it gives
$\sep_F(x)\leq\langle\,1\ /\ C\,\rangle\lVert f\otimes x\rVert_{X^*\projmodtens{A}F}$.
In $X^*\projmodtens{A}F$ the balancing relation
$(g\cdot a)\otimes x=g\otimes(a\cdot x)$ holds, so
$f\otimes x= \bigl(f-\sum_i g_i\cdot a_i \bigr)\otimes x
+\sum_i g_i\otimes(a_i\cdot x)$, whence the displayed bound by the triangle
inequality and $\lVert u\otimes v\rVert\leq\lVert u\rVert\lVert v\rVert$.
\end{proof}

\begin{corollary}\label{SepInFlatMod}
Let $F$ be a $\langle$~metrically / $C$-topologically~$\rangle$ flat left
$A$-module. Then for every $x\in F$, every $a\in A\setminus \{0 \}$ and every
$0<\gamma\leq 1$,
\[
\sep_F(x)\leq\langle\,1\ /\ C\,\rangle \Bigl(
  \max \{\gamma,1-\gamma \}\lVert x\rVert
  +\gamma\frac{\lVert a\cdot x\rVert}{\lVert a\rVert}
 \Bigr);
\]
with $\gamma=\tfrac{1}{2}$,
$\ 2\sep_F(x)\leq\langle\,1\ /\ C\,\rangle
\bigl(\lVert x\rVert+\lVert a\cdot x\rVert/\lVert a\rVert \bigr)$.
\end{corollary}

\begin{proof}
The statement quantifies over $a\in A\setminus \{0 \}$, so it is vacuous for
$A= \{0 \}$ and we may assume $A\neq \{0 \}$;
lemma~\ref{AugmCharDistEstimate} is then available. Apply
proposition~\ref{SepViaAnnihVector}$(iii)$ with $(X,f)=(A_+,\varepsilon)$ ---
admissible by lemma~\ref{AnnihPartOfUnitization} --- and with the single term
$g_1:=(\gamma/\lVert a\rVert)\varphi$, $a_1:=a$, where $\varphi\in A_+^*$
satisfies $\lVert\varphi\rVert=1$ and $\varphi(a/\lVert a\rVert)=1$. Then
$g_1\cdot a_1=\gamma\,\varphi\cdot(a/\lVert a\rVert)$, so
lemma~\ref{AugmCharDistEstimate} applied to the norm-one element
$a/\lVert a\rVert$ bounds
$\lVert\varepsilon-g_1\cdot a_1\rVert$ by $\max \{\gamma,1-\gamma \}$, while
$\lVert g_1\rVert\lVert a_1\cdot x\rVert
=\gamma\lVert a\cdot x\rVert/\lVert a\rVert$. Substituting the two bounds gives
the display, and $\max \{\gamma,1-\gamma \}=\tfrac{1}{2}$ at
$\gamma=\tfrac{1}{2}$.
\end{proof}

Taking the infimum over $a\neq 0$ in the second display gives the form in which
the estimate holds at every vector of $F$, not only the annihilator ones:
\[
2\sep_F(x)\leq\langle\,1\ /\ C\,\rangle \bigl(\lVert x\rVert+\sng(x) \bigr)
\qquad(x\in F),
\]
and $\gamma=\tfrac{1}{2}$ is the optimal choice there because
$\sng(x)\leq\lVert x\rVert$ always. This is not subsumed by the annihilator
estimate that follows, nor does it subsume it: the constant $\tfrac{1}{2}$ is the
value of $\sepd(A_+^*)$ only when $A$ has a contractive right approximate
identity.

\begin{corollary}\label{SepOfAnnihVectorInFlatMod}
Let $F$ be $\langle$~metrically / $C$-topologically~$\rangle$ flat and
$x\in F_{ann}$. Then
$\sep_F(x)\leq\langle\,1\ /\ C\,\rangle\sep_{X^*}(f)\lVert x\rVert$ for every
admissible pair $(X,f)$, that is, for every left $A$-module $X$ and every
$f\in{(X_{ess})}^\perp$ with $\lVert f\rVert=1$; and the best choice is
$(A_+,\varepsilon)$:
\[
\sep_F(x)\leq\langle\,\sepd(A_+^*)\ /\ C\sepd(A_+^*)\,\rangle\lVert x\rVert .
\]
\end{corollary}

\begin{proof}
For $x\in F_{ann}$ every term $\lVert a_i\cdot x\rVert$ in
proposition~\ref{SepViaAnnihVector}$(iii)$ vanishes, so that proposition reads
$\sep_F(x)\leq\langle\,1\ /\ C\,\rangle
\lVert f-\sum_i g_i\cdot a_i\rVert\lVert x\rVert$ for all finite families. The
elements $\sum_i g_i\cdot a_i$ are exactly the members of
$\operatorname{span} \{g\cdot a \}$, which is dense in ${(X^*)}_{ess}$; taking
the infimum over them replaces the norm by
$\dist \bigl(f,{(X^*)}_{ess} \bigr)=\sep_{X^*}(f)$.

For the best choice, theorem~\ref{SepOfDualUnitization}$(i)$ gives
$\sep_{X^*}(f)\geq\lVert f\rVert/(1+\crai(A))=\sepd(A_+^*)$ for every admissible
pair, so none yields a coefficient smaller than $\sepd(A_+^*)$; and the pair
$(A_+,\varepsilon)$ attains it, since
$\sep_{A_+^*}(\varepsilon)=\sepd(A_+^*)$ by
lemma~\ref{AnnihPartOfUnitization}. No hypothesis on $A$ is used.
\end{proof}

\begin{corollary}\label{FlatAnnihModHaveApprIdent}
Let $F\neq \{0 \}$ be an annihilator
$\langle$~metrically / $C$-topologically~$\rangle$ flat $A$-module. Then
$\langle$~$A= \{0 \}$ / $A$ has a right $(C-1)$-bounded approximate
identity~$\rangle$.
\end{corollary}

\begin{proof}
Since $F$ is a non-zero annihilator module, $F_{ess}= \{0 \}$ and $F_{ann}=F$, so
$\sep_F(x)=\lVert x\rVert$ for every $x\in F$; applying
corollary~\ref{SepOfAnnihVectorInFlatMod} at some $x\neq 0$ gives
$\sepd(A_+^*)\geq\langle\,1\ /\ 1/C\,\rangle$. In the metric bracket this says
$\sepd(A_+^*)\geq 1$, contradicting the ceiling $\tfrac{1}{2}$ of
lemma~\ref{AugmCharDistEstimate} unless the hypothesis of that lemma fails, that
is unless $A= \{0 \}$. In the topological bracket
theorem~\ref{SepOfDualUnitization}$(ii)$ turns it into $\crai(A)\leq C-1$, and
$(iii)$ of the same theorem says the bound is attained.
\end{proof}

\begin{corollary}\label{ExactSepOfAnnihVectorInFlatMod}
Let $F$ be a $\langle$~metrically / $C$-topologically~$\rangle$ flat left
$A$-module. Then, with no hypothesis on $A$,
\[
\frac{\lVert x\rVert}{1+\clai(A)}\ \leq\ \sep_F(x)\ \leq\
\langle\,\sepd(A_+^*)\ /\ C\sepd(A_+^*)\,\rangle\lVert x\rVert
\qquad(x\in F_{ann}).
\]
If moreover $\clai(A)\leq\crai(A)$, the metric bracket collapses to the equality
\[
\sep_F(x)=\sepd(A_+^*)\lVert x\rVert=\frac{\lVert x\rVert}{1+\crai(A)}
\qquad(x\in F_{ann}),
\]
which holds in particular when $A$ is commutative, when $A$ is unital, where the
value is $\lVert x\rVert/(1+\lVert e_A\rVert)$, and when $A\neq \{0 \}$ has a
contractive two-sided approximate identity, where it is $\tfrac{1}{2}\lVert
x\rVert$. Conversely, if some metrically flat left $A$-module has a non-zero
annihilator vector, then $\crai(A)\leq\clai(A)$.
\end{corollary}

\begin{proof}
The upper bound is corollary~\ref{SepOfAnnihVectorInFlatMod}. For the lower
bound, if $A$ has a left approximate identity then for every $\epsilon>0$ there
is one bounded by $K=\clai(A)+\epsilon$, and lemma~\ref{SepLowerBound} applied to
the left module $F$ gives $\sep_F(x)\geq\lVert x\rVert/(1+K)$; let
$\epsilon\to 0$. If $A$ has none then $\clai(A)=+\infty$ and the bound is vacuous.

Suppose $\clai(A)\leq\crai(A)$. Then
$\lVert x\rVert/(1+\clai(A))\geq\lVert x\rVert/(1+\crai(A))
=\sepd(A_+^*)\lVert x\rVert$ by
theorem~\ref{SepOfDualUnitization}$(ii)$, so in the metric bracket the lower
bound is at least the upper one and the two meet.

The three instances are the elementary relations recorded in
subsection~\ref{SubSectionBanachAlgebrasAndTheirModules}: a commutative algebra
has $\clai(A)=\crai(A)$; a unital one has
$\clai(A)=\crai(A)=\lVert e_A\rVert$, the zero algebra included, with $e_A=0$ and
value $\lVert x\rVert$; and an algebra $A\neq \{0 \}$ with a contractive
two-sided approximate identity has $\clai(A)=\crai(A)=1$, since such a net is
bounded by $1$ on both sides while every approximate identity of a non-zero
algebra is bounded below by $1$. The hypothesis $A\neq \{0 \}$ is not removable
in the third instance: the zero algebra has the zero net as a contractive
two-sided approximate identity, yet $\clai=\crai=0$ there and the value is
$\lVert x\rVert$.

For the converse, let $F$ be metrically flat with some $x\in F_{ann}$,
$x\neq 0$. Reading the displayed sandwich at that single $x$ and dividing by
$\lVert x\rVert$ gives $1/(1+\clai(A))\leq\sepd(A_+^*)=1/(1+\crai(A))$, that is
$\crai(A)\leq\clai(A)$.
\end{proof}

On the asymptotic layer the estimate simplifies: the master inequality bounds
$\sep_F$ at every vector, not only the annihilator ones, so it survives the
passage to nets.

\begin{proposition}\label{AsympSepOfFlatMod}
Let $A\neq \{0 \}$ and let $X$ be a
$\langle$~metrically / $C$-topologically~$\rangle$ flat left $A$-module. Then
$\asepd(X)\leq\langle\,1/2\ /\ C/2\,\rangle$; equivalently, every topological
flatness constant of $X$ satisfies $C\geq 2\asepd(X)$.
\end{proposition}

\begin{proof}
Since $A\neq \{0 \}$ we may fix $a_0\in A$ with $\lVert a_0\rVert=1$. Let
${(x_\nu)}_{\nu\in N}$ be an annihilator net in $B_X$;
corollary~\ref{SepInFlatMod} at $\gamma=\tfrac{1}{2}$ and $a=a_0$ gives
$2\sep_X(x_\nu)\leq\langle\,1\ /\ C\,\rangle
\bigl(1+\lVert a_0\cdot x_\nu\rVert \bigr)$. Take the limit superior, use
$\lVert a_0\cdot x_\nu\rVert\to 0$, and take the supremum over annihilator nets
in $B_X$.
\end{proof}

The metric bracket is sharp on free modules, which are metrically flat by
propositions~\ref{MetCTopFreeMod} and~\ref{MetTopProjIsMetTopFlat}: by
proposition~\ref{AsympSepOfFreeMod} they have $\asepd=\sepd(A_+^*)$, and that
equals $\tfrac{1}{2}$ exactly when $A$ has a contractive right approximate
identity.

\begin{corollary}\label{AsympSepQuantization}
Let $A\neq \{0 \}$ have a right approximate identity. Then every non-essential
$\langle$~metrically / $C$-topologically~$\rangle$ flat left $A$-module $X$
satisfies
\[
\frac{1}{1+\crai(A)}\leq\asepd(X)\leq\langle\,1/2\ /\ C/2\,\rangle .
\]
In particular, over an algebra with a contractive right approximate identity the
asymptotic separation defect of a metrically flat module takes exactly two
values: $\asepd(X)=0$ if $X$ is essential and $\asepd(X)=\tfrac{1}{2}$ otherwise.
\end{corollary}

\begin{proof}
The upper bound is proposition~\ref{AsympSepOfFlatMod}. For the lower one, $X$ is
not essential, so proposition~\ref{TheChain} and
theorem~\ref{SepOfDualUnitization}$(i)$, $(ii)$ give
$\asepd(X)\geq\sepd(X^*)\geq\sepd(A_+^*)=1/(1+\crai(A))$. If $A$ has a
contractive right approximate identity then $\crai(A)=1$ by the first clause of
corollary~\ref{SepOfDualUnitizationValues}, so for non-essential metrically flat
$X$ the display reads $\tfrac{1}{2}\leq\asepd(X)\leq\tfrac{1}{2}$; and if $X$ is
essential then $\sep_X$ vanishes identically.
\end{proof}

\begin{corollary}\label{MaxAsympSepImpliesNotFlat}
Let $A\neq \{0 \}$. A left $A$-module of asymptotic separation defect $1$ is
never metrically flat, and its topological flatness constants satisfy $C\geq 2$.
\end{corollary}

\begin{proof}
If $X$ were metrically flat, proposition~\ref{AsympSepOfFlatMod} would give
$\asepd(X)\leq\tfrac{1}{2}$, contradicting $\asepd(X)=1$; and if $X$ is
$C$-topologically flat the same proposition gives $1\leq C/2$.

The hypothesis $A\neq \{0 \}$ is needed, and $\mathbb{C}$ is the witness. Over
the zero algebra every module is an annihilator module, so
$\mathbb{C}_{ess}= \{0 \}$, $\sep_{\mathbb{C}}=\lVert\cdot\rVert$ and
$\asepd(\mathbb{C})=1$; on the other hand, as recorded in
subsection~\ref{SubSectionMetricAndTopologicalFlatness}, modules over the zero
algebra are Banach spaces and metric flatness for them is Grothendieck's
$L_1$ condition \cite[theorem 1]{GrothMetrProjFlatBanSp}, which $\mathbb{C}$
satisfies as the $L_1$-space of a one-point measure space. So $\mathbb{C}$ is a
metrically flat module of asymptotic separation defect $1$ over $A= \{0 \}$.
\end{proof}

\subsection{Homologically trivial annihilator modules}\label{SubSectionHomoligicallyTrivialAnnihilatorModules}

\emph{This subsection is the present subsection 2.2.1 of [B], moved here
unchanged in content and order:} proposition~\ref{AnnihCModIsRetAnnihMod},
proposition~\ref{MetTopProjModCCharac},
proposition~\ref{MetTopProjOfAnnihModCharac},
proposition~\ref{MetTopInjModCCharac},
proposition~\ref{MetTopInjOfAnnihModCharac},
proposition~\ref{MetTopFlatAnnihModCharac}, and the proposition asserting that no
Banach algebra makes all of its modules flat, which is unlabelled in [B] and
needs a label before it can be cited (\texttt{NoAlgWithAllModFlat} is proposed in
Part~\ref{PartIV}). \emph{What changes is the proof of the algebra clause of each
$(i)\Rightarrow(ii)$ half, and two corollaries are added.} The three
characterisations then have one proof with two constants: $\sepd(A_+)$ for the
projective case and $\sepd(A_+^*)$ for the other two.

\paragraph{The algebra clause of
proposition~\ref{MetTopProjModCCharac}$(i)\Rightarrow(ii)$ and of
proposition~\ref{MetTopProjOfAnnihModCharac}$(i)\Rightarrow(ii)$.} Let $P$ be a
non-zero annihilator $A$-module which is
$\langle$~metrically / $C$-topologically~$\rangle$ projective; the case
$P=\mathbb{C}$ is proposition~\ref{MetTopProjModCCharac}. By
corollary~\ref{ProjAnnihMod}, $\langle$~$A= \{0 \}$ / $\cri(A)\leq C-1$~$\rangle$.
This is one application of theorem~\ref{SepOfProjMod} to the single lifting
problem $\varepsilon:A_+\to\mathbb{C}$, in place of the idempotent argument of
proposition~\ref{MetTopProjCycModCharac}, and by
corollary~\ref{ProjOverNonUnital} it says more than the existence of the
identity: over an algebra without a right identity every annihilator vector of a
projective module already lies in its essential part.

\emph{One caveat on the topological bracket.} The clause as printed in [B] asks
for a right identity of norm at most $C-1$, and the separation estimate gives
$\cri(A)\leq C-1$, that is, a right identity of norm at most $C-1+\epsilon$ for
every $\epsilon>0$. The two are not the same, because the infimum defining
$\cri(A)$ need not be attained: the right identities of $A$ form a closed coset
$u+ \{z\in A:Az= \{0 \} \}$, and closed subspaces of a Banach space are not in
general proximinal. This is exactly where the projective constant differs from
the dual one, whose infimum \emph{is} attained
(theorem~\ref{SepOfDualUnitization}$(iii)$). Either restate the clause as
$\cri(A)\leq C-1$ --- which is the invariant form, and which we recommend --- or
keep the attained form and with it the appeal to
proposition~\ref{MetTopProjCycModCharac}, which is legal here since that
proposition lies in section 2.1.

\paragraph{The algebra clause of
proposition~\ref{MetTopInjModCCharac} and of
proposition~\ref{MetTopInjOfAnnihModCharac}$(i)\Rightarrow(ii)$.} Let
$J\neq \{0 \}$ be an annihilator right $A$-module which is
$\langle$~metrically / $C$-topologically~$\rangle$ injective. By
corollary~\ref{InjAnnihModApprIdent}, $\langle$~$A= \{0 \}$ / $A$ has a right
$(C-1)$-bounded approximate identity~$\rangle$; here the bound \emph{is}
attained, by theorem~\ref{SepOfDualUnitization}$(iii)$. This is three lines and
replaces the passage through metric and topological flatness of $A_+/A$.

\paragraph{The algebra clause of
proposition~\ref{MetTopFlatAnnihModCharac}$(i)\Rightarrow(ii)$.} Identical, with
corollary~\ref{FlatAnnihModHaveApprIdent} in place of
corollary~\ref{InjAnnihModApprIdent}.

The constant $C-1$ is moreover optimal, and the least admissible $C$ is the
reciprocal of the dual constant.

\begin{corollary}\label{InjAnnihModApprIdentOptimal}
Let $J\neq \{0 \}$ be an annihilator right $A$-module which is metrically
injective as a Banach space. Then
\[
\inf \{C:J\ \text{is a $C$-topologically injective $A$-module} \}
=\frac{1}{\sepd(A_+^*)}=1+\crai(A),
\]
and the infimum is attained; it is infinite exactly when $A$ has no right
approximate identity. The same holds, with the same proof, for an annihilator
$A$-module $F\neq \{0 \}$ that is metrically flat as a Banach space and for the
constant of proposition~\ref{MetTopFlatAnnihModCharac}.
\end{corollary}

\begin{proof}
By corollary~\ref{InjAnnihModApprIdent} every admissible $C$ satisfies
$C\geq 1+\crai(A)=1/\sepd(A_+^*)$, which is also the assertion about
finiteness. For the upper bound let $\crai(A)<+\infty$ and put
$C:=1+\crai(A)$. By theorem~\ref{SepOfDualUnitization}$(iii)$ the algebra $A$ has
a right $(C-1)$-bounded approximate identity; and $J$, being metrically injective
as a Banach space, is $1$-topologically injective as a Banach space by
proposition~\ref{MetInjIsOneTopInj} read over the zero algebra, hence
$C$-topologically injective as a Banach space for every $C\geq 1$. These are
exactly the hypotheses of
proposition~\ref{MetTopInjOfAnnihModCharac} in the direction
$(ii)\Rightarrow(i)$, whose conclusion is that $J$ is $C$-topologically injective
as an $A$-module. The flat statement is the same argument with
proposition~\ref{MetTopFlatAnnihModCharac} in place of
proposition~\ref{MetTopInjOfAnnihModCharac}, a metrically flat Banach space being
an $\mathscr{L}_{1,1}^g$-space and hence an $\mathscr{L}_{1,C}^g$-space for every
$C\geq 1$.
\end{proof}

Finally, the asymptotic layer carries the same three characterisations one step
further: propositions~\ref{AsympSepOfProjMod} and~\ref{AsympSepOfFlatMod} with
corollaries~\ref{AsympSepQuantization} and~\ref{MaxAsympSepImpliesNotFlat} say
that over an algebra with a contractive right approximate identity a metrically
flat module has $\asepd\in \{0,\tfrac{1}{2} \}$, and that a module of asymptotic
separation defect $1$ over a non-zero algebra is never metrically flat. That last
statement is what the measure algebra needs, and it is proved in
Part~\ref{PartIII}.

\section{For chapter 3}\label{PartIII}

\subsection{For subsection~\ref{SectionPreliminariesOnHarmonicAnalysis}, replacing
the sentences on the Lebesgue decomposition of $M(G)$}

\emph{Delete} the assertions that $M(G)$ is a coproduct of $M_a(G)$ and $M_s(G)$
in $L_1(G)\text{-}\mathbf{mod}_1$, that $M_s(G)$ is a subalgebra of $M(G)$, and
that $M_s(G)$ is an annihilator $L_1(G)$-module. All three fail for non-discrete
$G$, because $M_s(G)$ is then not even an $L_1(G)$-submodule of $M(G)$: points
are $m_G$-null, so $\delta_x\in M_s(G)$ for every $x\in G$, while
$gm_G\convol\delta_x$ is non-zero and absolutely continuous for every $g\neq 0$
in $L_1(G)$; hence $L_1(G)\convol M_s(G)\not\subset M_s(G)$ and
$M_s(G)\convol L_1(G)\neq \{0 \}$. Nor is $M_s(G)$ a subalgebra: Wiener and
Wintner constructed on $\mathbb{T}$ a singular probability measure $\mu$ with
$\widehat\mu(n)=O(|n|^{-1/2+\epsilon})$, and then $\mu^{\convol 4}$ has
absolutely convergent Fourier series and is a non-zero absolutely continuous
measure.

\emph{Insert} instead: $M_a(G)$ is a closed two-sided ideal of $M(G)$ with
$M_a(G)\isom{M(G)\text{-}\mathbf{mod}_1}L_1(G)$; the Lebesgue decomposition
$M(G)\isom{\mathbf{Ban}_1}M_a(G)\bigoplus_1 M_s(G)$ is an isomorphism of Banach
spaces and of Banach spaces only; and ${M(G)}_{ess}=M_a(G)$, so that
$M(G)/{M(G)}_{ess}\isom{\mathbf{Ban}_1}M_s(G)$, this last being proved in
proposition~\ref{AsympSepOfMeasAlg} below. That is what the proofs of
proposition~\ref{L1MetTopProjAndMetrFlatOfMeasAlg}$(iii)$ and of
proposition~\ref{MeasAlgIsL1TopFlat} already use, and all that they need; the
opening paragraph of the latter, which proves ${M(G)}_{ess}=M_a(G)$ a second
time, can be cut and replaced by a citation.

One sentence is worth adding at the point of deletion, because it is how the
error arose. Subsection~\ref{SubSectionC0SModules} correctly asserts that
$M_s(S,\mu)$ \emph{is} a closed $C_0(S)$-submodule of $M(S)$: there the action is
pointwise multiplication, which preserves singularity with respect to $\mu$.
Here the action is convolution, which does not.

\subsection{For subsection~\ref{SubSectionL1GModules}}

The following proposition goes at the head of the subsection. It is the
supporting statement for the correction above and for the replacement proof
that follows it.

\begin{proposition}\label{AsympSepOfMeasAlg}
Let $G$ be a non-discrete locally compact group and regard $M(G)$ as a left
$L_1(G)$-module. Then ${M(G)}_{ess}=M_a(G)$ and ${M(G)}_{ann}= \{0 \}$, and
\[
\asepd(M(G))=1,\qquad \sepd(M{(G)}^*)=1,\qquad \sepd(M(G))=0 .
\]
\end{proposition}

\begin{proof}
Points are Haar-null. By left invariance $m_G( \{x \})=m_G( \{e_G \})=:c$ for
every $x\in G$. If $c>0$, take a compact neighbourhood $W$ of $e_G$, so that
$m_G(W)<+\infty$; any $n$ distinct points of $W$ give $m_G(W)\geq nc$, so $W$ is
finite, and a finite neighbourhood of $e_G$ in a Hausdorff group contains a
finite open set, which makes $ \{e_G \}$ open and $G$ discrete. Hence $c=0$ and
$\delta_x\in M_s(G)$ for every $x\in G$.

The essential part. On the one hand $L_1(G)\convol M(G)\subset M_a(G)$, a
convolution with an absolutely continuous measure being absolutely continuous; on
the other hand $M_a(G)$ is isometrically isomorphic to $L_1(G)$ as a module and
$L_1(G)$ has a contractive approximate identity, so $M_a(G)$ is essential. Hence
${M(G)}_{ess}=M_a(G)$. The annihilator part is zero: if $f\convol\mu=0$ for every
$f\in L_1(G)$, take $f=\chi_U/m_G(U)$ over a shrinking neighbourhood base at
$e_G$, so that $f\convol\mu\to\mu$ weak${}^*$ and $\mu=0$. Consequently
$B_{{M(G)}_{ann}}= \{0 \}$ and $\sepd(M(G))=0$: the measure algebra is a witness
for the strictness of the first inequality of proposition~\ref{TheChain}, and it
realises the maximal asymptotic defect without a single annihilator vector.

The separation of a difference of point masses. For $z\neq e_G$ put
$\mu_z:=\tfrac{1}{2}(\delta_{e_G}-\delta_z)$. The two point masses are mutually
singular, so $\lVert\mu_z\rVert=1$, and $\mu_z\in M_s(G)$; the
$\bigoplus_1$-decomposition gives
$\lVert\mu_z-\nu\rVert=\lVert\mu_z\rVert+\lVert\nu\rVert\geq 1$ for every
$\nu\in M_a(G)$, whence $\sep_{M(G)}(\mu_z)=1$.

An annihilator net at separation one. Let $N$ be a neighbourhood base at $e_G$,
directed by reverse inclusion. Since $G$ is not discrete, no $V\in N$ reduces to
$ \{e_G \}$, so we may choose $z_V\in V\setminus \{e_G \}$. For $g\in L_1(G)$ one
has $g\convol\mu_{z_V}=\tfrac{1}{2}(g-g\convol\delta_{z_V})$, and
$\lVert g-g\convol\delta_z\rVert_1\to 0$ as $z\to e_G$ by the norm continuity of
translation in $L_1(G)$ together with the continuity of the modular function; so
${(\mu_{z_V})}_{V\in N}$ is an annihilator net in $B_{M(G)}$ with
$\sep_{M(G)}(\mu_{z_V})=1$ at every index, and $\asepd(M(G))\geq 1$. As
$\asepd\leq 1$ always, this is an equality.

The dual. Every $\mu\in M(G)$ has at most countably many atoms and
$\sum_x|\mu( \{x \})|\leq\lVert\mu\rVert$, so
$f(\mu):=\mu( \{e_G \})-\sum_{x\neq e_G}\mu( \{x \})$ is a well-defined linear
functional with $\lVert f\rVert\leq 1$. It vanishes on $M_a(G)$, an absolutely
continuous measure having no atoms, that is
$f\in{({M(G)}_{ess})}^\perp={\bigl(M{(G)}^* \bigr)}_{ann}$; and
$f(\mu_z)=\tfrac{1}{2}+\tfrac{1}{2}=1$, so $\lVert f\rVert=1$ and
proposition~\ref{SepOfDualViaAnnihNets} applied to the net above gives
$\sep_{M{(G)}^*}(f)\geq 1$. Hence $\sepd(M{(G)}^*)=1$.
\end{proof}

\emph{Replacing the proof of
proposition~\ref{L1MetTopProjAndMetrFlatOfMeasAlg}$(iv)\Rightarrow(i)$.} Delete
the present argument, which opens with
$M(G)\isom{L_1(G)\text{-}\mathbf{mod}_1}L_1(G)\bigoplus_1 M_s(G)$ and concludes
through propositions~\ref{MetTopFlatModCoProd} and~\ref{MetTopFlatAnnihModCharac}.
Insert:

\begin{quote}
$(iv)\Rightarrow(i)$ Suppose $G$ is not discrete. By
proposition~\ref{AsympSepOfMeasAlg} the $L_1(G)$-module $M(G)$ has asymptotic
separation defect $1$, and $L_1(G)\neq \{0 \}$; so
corollary~\ref{MaxAsympSepImpliesNotFlat} says that $M(G)$ is not metrically
flat. Hence a metrically flat $M(G)$ forces $G$ to be discrete.
\end{quote}

Items $(i)$--$(iii)$ of that proposition stand as printed. The same estimate
bounds the topological constant from below, which the present text does not
record.

\begin{corollary}\label{MeasAlgL1TopFlatConst}
Let $G$ be a non-discrete locally compact group. Then every constant of
topological flatness of the $L_1(G)$-module $M(G)$ is at least $2$. The module is
topologically flat by proposition~\ref{MeasAlgIsL1TopFlat}, so the least such
constant lies in $[2,+\infty)$.
\end{corollary}

\begin{proof}
By proposition~\ref{AsympSepOfMeasAlg} we have $\asepd(M(G))=1$, and
$L_1(G)\neq \{0 \}$; apply corollary~\ref{MaxAsympSepImpliesNotFlat}.
\end{proof}

\subsection{For subsection~\ref{SubSectionc0AndlInftyModules}, the metric half of
proposition~\ref{c0AndlInftyModlIfty}$(ii)$}

\emph{Replace} the appeal to proposition~\ref{CStarAlgIsTopFlatOverItsIdeal} by
the asymptotic separation defect. The paragraph becomes:

\begin{quote}
$(ii)$ Let $\Lambda$ be infinite. We claim that
$\asepd(\ell_\infty(\Lambda))=1$ over $c_0(\Lambda)$. The essential part is
$\operatorname{cl}(c_0(\Lambda)\ell_\infty(\Lambda))=c_0(\Lambda)$. For a finite
$S\subset\Lambda$ put $v_S:=\chi_{\Lambda\setminus S}$; these are unit vectors,
they satisfy $\lVert av_S\rVert_\infty=\sup_{\lambda\notin S}|a(\lambda)|\to 0$
for every $a\in c_0(\Lambda)$ as $S$ runs over the finite subsets of $\Lambda$
directed by inclusion, so they form an annihilator net in
$B_{\ell_\infty(\Lambda)}$; and
$\sep_{\ell_\infty(\Lambda)}(v_S)=\dist(v_S,c_0(\Lambda))=1$, a norming direction
being $f(x)=\lim_{\mathfrak{U}}x(\lambda_n)$ along a free ultrafilter
$\mathfrak{U}$ on $\mathbb{N}$ and a sequence of distinct
$\lambda_n\in\Lambda\setminus S$, which vanishes on $c_0(\Lambda)$, has norm at
most one and takes the value $1$ at $v_S$. Since $c_0(\Lambda)\neq \{0 \}$,
corollary~\ref{MaxAsympSepImpliesNotFlat} applies and gives that
$\ell_\infty(\Lambda)$ is not a metrically flat $c_0(\Lambda)$-module. Both
claims regarding projectivity follow from
proposition~\ref{MetTopProjIsMetTopFlat}. If $\Lambda$ is finite, then
$c_0(\Lambda)=\ell_\infty(\Lambda)$ and the result follows from paragraph $(i)$.
\end{quote}

The same estimate gives, as prose after the proposition, that every constant of
topological flatness of the $c_0(\Lambda)$-module $\ell_\infty(\Lambda)$ is at
least $2$ for infinite $\Lambda$ --- a quantitative refinement that the appeal to
proposition~\ref{CStarAlgIsTopFlatOverItsIdeal} does not give. The topological
half of $(ii)$ is unaffected and keeps its present proof.

\section{Merge manifest}\label{PartIV}

\subsection{Preamble additions required in \texttt{book/main.tex}}

The following go beside the existing definitions of \texttt{\textbackslash
projtens}, \texttt{\textbackslash convol}, \texttt{\textbackslash projmodtens}
and \texttt{\textbackslash isom}. Nothing else in the preamble changes, and in
particular \texttt{\textbackslash sp} is left alone: [B] uses
$\operatorname{sp}_A(a)$ for the spectrum of an algebra element throughout
chapter 3, so the separation operator is called $\sep$ and not $\operatorname{sp}$.

\begin{verbatim}
\DeclareMathOperator{\dist}{dist}
\DeclareMathOperator{\sep}{sep}
\DeclareMathOperator{\sepd}{sepd}
\DeclareMathOperator{\asep}{asep}
\DeclareMathOperator{\asepd}{asepd}
\DeclareMathOperator{\sng}{sng}
\newcommand{\cri}{c_{\mathrm{ri}}}
\newcommand{\cli}{c_{\mathrm{li}}}
\newcommand{\crai}{c_{\mathrm{rai}}}
\newcommand{\clai}{c_{\mathrm{lai}}}
\end{verbatim}

The four constants may equally be spelled out in the text; the operators cannot,
and $\sng$ is the only one of the six whose name was free in [B] already.

\subsection{Chapter 1}

\begin{itemize}
\item \emph{Insertion}, subsection~\ref{SubSectionBanachSpaces}, in the paragraph
carrying $E^*/F^\perp\isom{\mathbf{Ban}_1}F^*$: the $\dist$ notation, the
Hahn-Banach distance formula, Banach-Alaoglu, Goldstine and Mazur. Part I.1.
Prose, no numbered statement.

\item \emph{Insertion}, subsection~\ref{SubSectionBanachAlgebrasAndTheirModules},
after the paragraph on approximate identities: the four constants
$\cri,\cli,\crai,\clai$, their conventions, the elementary relations, the open
question, and the forward pointer to attainment. Part I.2.

\item \emph{Insertion}, same subsection: the definition of an annihilator net,
with the disclaimer that no norm bound is part of it. Part I.3.

\item \emph{Insertion}, same subsection, in the paragraph that already records
${(X^*)}^{\perp S}={\operatorname{cl}_X(SX)}^\perp$: the two dual identities
spelled out at $S=A$. Part I.4. Note that [B] already has the first of them; only
the second, and the remark that ${(X^*)}_{ess}\subset{(X_{ann})}^\perp$ can be
proper, are new.
\end{itemize}

\subsection{Chapter 2}

\begin{itemize}
\item \emph{Renumbering.} A new section 2.2, ``The essential and the annihilator
part'', is inserted between the present 2.1 and the present 2.2. The present
subsection 2.2.1, \texttt{SubSectionHomoligicallyTrivialAnnihilatorModules},
becomes subsection 2.2.6 of the new section; the rest of the present section 2.2
becomes section 2.3, and the present section 2.3 becomes section 2.4. No label of
[B] changes, and no statement of [B] is renumbered by hand: every citation in
this note is by label.

\item \emph{Relocation, no change of content}, into
subsection~\ref{SubSectionAnnihPartOfProjMod} at its head:
proposition~\ref{MetTopRelProjModCompIdealPartCompl},
corollary~\ref{MetTopRelProjModPartCompl},
proposition~\ref{MetTopProjModEssQuotIsProjBanSp}.

\item \emph{Relocation, no change of content}, into
subsection~\ref{SubSectionAnnihPartOfInjMod} at its head:
proposition~\ref{MetTopRelInjModCompIdealAnnihCompl},
corollary~\ref{MetTopRelInjModAnnihCompl},
proposition~\ref{MetTopInjModAnnihPartIsInjBanSp}.

\item \emph{Relocation, no change of content}, into
subsection~\ref{SubSectionAnnihPartOfFlatMod} at its head:
proposition~\ref{MetTopRelFlatModCompIdealPartCompl},
corollary~\ref{MetTopRelFlatModPartCompl},
proposition~\ref{MetTopFlatModEssQuotIsFlatBanSp}.

\item \emph{Insertion}, subsections~\ref{SubSectionSeparation}
to~\ref{SubSectionAnnihPartOfFlatMod}: everything else in Part II, in the order
printed there.

\item \emph{Replacement}, proposition~\ref{MetTopProjModCCharac} and
proposition~\ref{MetTopProjOfAnnihModCharac}, the algebra clause of
$(i)\Rightarrow(ii)$: by corollary~\ref{ProjAnnihMod}. See the caveat in
subsection~\ref{SubSectionHomoligicallyTrivialAnnihilatorModules} and item
(c) below.

\item \emph{Replacement}, proposition~\ref{MetTopInjModCCharac} and
proposition~\ref{MetTopInjOfAnnihModCharac}, the algebra clause of
$(i)\Rightarrow(ii)$: by corollary~\ref{InjAnnihModApprIdent}.

\item \emph{Replacement}, proposition~\ref{MetTopFlatAnnihModCharac}, the algebra
clause of $(i)\Rightarrow(ii)$: by corollary~\ref{FlatAnnihModHaveApprIdent}.

\item \emph{Insertion}, subsection
\ref{SubSectionHomoligicallyTrivialAnnihilatorModules}, after
proposition~\ref{MetTopFlatAnnihModCharac}:
corollary~\ref{InjAnnihModApprIdentOptimal}.

\item \emph{Labelling}, subsection
\ref{SubSectionHomoligicallyTrivialAnnihilatorModules}: the proposition
asserting that no Banach algebra makes all of its modules flat carries no label
in [B] and cannot be cited. Proposed label:
\texttt{NoAlgWithAllModFlat}.
\end{itemize}

\subsection{Chapter 3}

\begin{itemize}
\item \emph{Deletion and replacement},
subsection~\ref{SectionPreliminariesOnHarmonicAnalysis}: the coproduct, the
subalgebra and the annihilator-module assertions about $M_s(G)$, replaced by the
text of Part III.1.

\item \emph{Insertion}, subsection~\ref{SubSectionL1GModules} at its head:
proposition~\ref{AsympSepOfMeasAlg}.

\item \emph{Replacement}, proposition~\ref{L1MetTopProjAndMetrFlatOfMeasAlg},
the proof of $(iv)\Rightarrow(i)$: by the three lines of Part III.2. Items
$(i)$--$(iii)$ stand as printed.

\item \emph{Insertion}, subsection~\ref{SubSectionL1GModules}, after
proposition~\ref{MeasAlgIsL1TopFlat}:
corollary~\ref{MeasAlgL1TopFlatConst}.

\item \emph{Deletion}, proposition~\ref{MeasAlgIsL1TopFlat}: its opening
paragraph, which proves ${M(G)}_{ess}=M_a(G)$ a second time, replaced by a
citation of proposition~\ref{AsympSepOfMeasAlg}.

\item \emph{Replacement}, proposition~\ref{c0AndlInftyModlIfty}$(ii)$, the metric
half: by the paragraph of Part III.3, which removes the appeal to
proposition~\ref{CStarAlgIsTopFlatOverItsIdeal}. The topological half is
unaffected.
\end{itemize}

\subsection{Points to vet before merging}

\begin{enumerate}[label = (\alph*)]
\item \emph{The $K=0$ gap in lemma~\ref{SepLowerBound}} is closed by disposing of
$K=0$ in a sentence, not by imposing $K\geq 1$. Imposing $K\geq 1$ would be
consistent with [B]'s notion of a $c$-bounded approximate identity, but it would
break the lemma's own consequence at $A= \{0 \}$, where
corollary~\ref{ExactSepOfAnnihVectorInFlatMod} invokes it with
$K=\clai(A)+\epsilon=\epsilon<1$. Since $K\geq 1$ is automatic for
$A\neq \{0 \}$, the hypothesis would bite only on the zero algebra --- the exact
case the rest of the sweep is trying to include.

\item \emph{The zero algebra in
subsection~\ref{SubSectionSeparationOfUnitization}.} The hypothesis
$A\neq \{0 \}$ has been dropped from theorem~\ref{SepOfDualUnitization}
altogether, and all three clauses were checked at $A= \{0 \}$:
${(A_+)}_{ess}= \{0 \}$, ${(A_+^*)}_{ann}=A_+^*=\mathbb{C}\varepsilon$,
${(A_+^*)}_{ess}= \{0 \}$, so $\sepd(A_+^*)=1=1/(1+\crai( \{0 \}))$; the zero net
is admissible in the lower bound of $(ii)$ and attains the value in $(iii)$,
where $\delta=1$ forces $m=0$ and $p=0$; and $(i)$ reads
$\sep_{X^*}(f)\geq\lVert f\rVert$, which holds with equality because
$X_{ess}= \{0 \}$. The hypothesis is kept exactly where it bites:
lemma~\ref{AugmCharDistEstimate}, which needs an element of norm one, and hence
corollary~\ref{SepOfDualUnitizationValues}, whose first clause fails at
$A= \{0 \}$. The asymmetry with lemma~\ref{AnnihPartOfUnitization}, which never
needed the hypothesis, is thereby removed, and the zero-algebra paragraph
formerly buried in the proof of
corollary~\ref{SepOfAnnihVectorInFlatMod} disappears, along with the case splits
it forced on corollaries~\ref{FlatAnnihModHaveApprIdent}
and~\ref{ExactSepOfAnnihVectorInFlatMod}.

\item \emph{Attainment of $\cri(A)$.} The separation estimate proves
$\cri(A)\leq C-1$, whereas propositions~\ref{MetTopProjModCCharac}
and~\ref{MetTopProjOfAnnihModCharac} as printed assert a right identity of norm
at most $C-1$. The infimum defining $\cri(A)$ need not be attained: the right
identities of $A$ form the closed coset $u+ \{z\in A:Az= \{0 \} \}$, and closed
subspaces of a Banach space are not in general proximinal. This is a real
difference from the injective and flat clauses, where
theorem~\ref{SepOfDualUnitization}$(iii)$ does give attainment. Recommendation:
restate the two projective clauses as $\cri(A)\leq C-1$. If the attained form is
to be kept, the appeal to proposition~\ref{MetTopProjCycModCharac} must be kept
with it; that is legal, since it lies in section 2.1.

\item \emph{Theorem~\ref{SepOfDualUnitization}$(iii)$ has been retained.} The
instruction was to check whether the attainment of $\crai(A)$ is classical ---
Doran and Wichmann, \cite{AppIdAndFactorInBanAlg} --- and, if so, to cite it and
delete $(iii)$. That could not be verified against the source, and deleting a
short self-contained proof on the strength of an unverified citation seemed the
worse risk. The clause is used in
corollaries~\ref{InjAnnihModApprIdent},~\ref{FlatAnnihModHaveApprIdent}
and~\ref{InjAnnihModApprIdentOptimal} and in
corollary~\ref{SepOfDualUnitizationValues}, so if the reference is confirmed it
is those four citations that change, not the argument.

\item \emph{One label is split.} The former
\texttt{InjAnnihModApprIdentOptimal} had two clauses. Clause $(i)$ is
corollary~\ref{InjAnnihModApprIdent} and stays in
subsection~\ref{SubSectionAnnihPartOfInjMod}; clause $(ii)$ keeps the old label
and moves to subsection
\ref{SubSectionHomoligicallyTrivialAnnihilatorModules}. The move is forced: the
proof of $(ii)$ quotes propositions~\ref{MetTopInjOfAnnihModCharac}
and~\ref{MetTopFlatAnnihModCharac} in the direction $(ii)\Rightarrow(i)$, and
those live in subsection 2.2.6. Leaving it in 2.2.4 would be a forward reference
of exactly the kind the reorganization is meant to remove.

\item \emph{One statement is shortened.} Corollary~\ref{ProjAnnihMod} is stated
as a necessary condition only. Its converse, in the note, was a citation of
proposition~\ref{MetTopProjOfAnnihModCharac}, which now sits later; and since the
corollary dissolves into that very proposition, nothing is lost.

\item \emph{Proof-internal cross-references removed, beyond the four the
instruction listed.} Three more were found and are handled the same way. The
identity $X_{ann}=\bigoplus_p \{{(X_\lambda)}_{ann} \}$ is now part of the
statement of proposition~\ref{SepOfSums}, since
corollary~\ref{SepOfInjMod}$(ii)$ used it. The identities
${(A_+^*)}_{ann}=\mathbb{C}\varepsilon$, $\lVert\varepsilon\rVert=1$ and
$\sepd(A_+^*)=\sep_{A_+^*}(\varepsilon)$ are now part of the statement of
lemma~\ref{AnnihPartOfUnitization}, since
proposition~\ref{SepOfAnnihVectorInInjMod},
corollary~\ref{SepOfAnnihVectorInFlatMod} and
proposition~\ref{AsympSepOfFreeMod} all quoted them from inside the proof of
theorem~\ref{SepOfDualUnitization}$(ii)$. And
proposition~\ref{SepOfCofreeMod}, which quoted the essential part of the free
module from inside the proof of proposition~\ref{AsympSepOfFreeMod}, now derives
it from proposition~\ref{SepOfSums}$(iv)$ directly.

\item \emph{The reading of proposition~\ref{MetTopFlatCycModCharac}} recorded at
the end of subsection~\ref{SubSectionSeparationOfUnitization} is stated for the
unitization $A_+$. It is there, and only there, that
$\lVert e_{A_+}-e_\nu\rVert=1+\lVert e_\nu\rVert$ for $e_\nu\in I\subset A$; over
$A_\times$ the quantity is not $1/\sepd(I_+^*)$.

\item \emph{Length.} No proof in Part II exceeds the longest proofs of chapter 2
of [B], namely those of lemma~\ref{GoodIdealMetTopProjIsUnital} and
proposition~\ref{MetTopFlatCycModCharac}. The step-numbered format is gone
everywhere, and the verifications [B] habitually compresses --- three-epsilon
arguments, checks that a constructed map is a morphism, directing a net over
pairs $(F,\epsilon)$, and the $1$-Lipschitz and positively homogeneous character
of $\sep_X$ --- are compressed here as well, the last being stated once after
definition~\ref{SepOfMod} and used silently thereafter.
\end{enumerate}

\end{document}
